\documentclass[11pt,letterpaper]{article}
\usepackage[T1]{fontenc}
\usepackage{lmodern,microtype,amsmath,amssymb,amsthm,mathtools}
\usepackage[margin=0.85in]{geometry}
\usepackage{booktabs,longtable,array,multicol,enumitem}
\usepackage{xcolor,xurl,hyperref,fancyhdr}
\hypersetup{colorlinks=true,linkcolor=blue!45!black,urlcolor=blue!45!black,citecolor=blue!45!black,
 pdftitle={Shared Pell bounds for modular certificates},pdfauthor={Research prepared for Vladimir Reshetnikov}}
\setlength{\parskip}{4pt}
\setlength{\parindent}{0pt}
\setlength{\emergencystretch}{2em}
\setlist{nosep,leftmargin=*}
\pagestyle{fancy}\fancyhf{}\fancyhead[L]{\small Diophantine consistency: shared index certificates}
\fancyhead[R]{\small 24 October 2026}\fancyfoot[C]{\thepage}
\setlength{\headheight}{14pt}
\newcommand{\N}{\mathbb N}
\newcommand{\Z}{\mathbb Z}
\newcommand{\Pos}{\mathbb Z_{>0}}
\newcommand{\Con}{\operatorname{Con}}
\newcommand{\ZFC}{\mathsf{ZFC}}
\newcommand{\ZF}{\mathsf{ZF}}
\newcommand{\degr}{\operatorname{deg}}
\newcommand{\Height}{\operatorname{ht}}
\newcommand{\floor}[1]{\left\lfloor#1\right\rfloor}
\newcommand{\ceil}[1]{\left\lceil#1\right\rceil}
\newcommand{\pathfile}[1]{\nolinkurl{#1}}
\newtheorem{theorem}{Theorem}[section]
\newtheorem{lemma}[theorem]{Lemma}
\newtheorem{proposition}[theorem]{Proposition}
\newtheorem{corollary}[theorem]{Corollary}
\theoremstyle{remark}\newtheorem{remark}[theorem]{Remark}
\newcommand{\MainVars}{241}
\newcommand{\MainRows}{189}
\newcommand{\MainTerms}{30,058}
\newcommand{\MainHeight}{85,187,728}
\newcommand{\DenseVars}{239}
\newcommand{\DenseRows}{187}
\newcommand{\DenseTerms}{387,494}
\newcommand{\DenseHeight}{269,805,528}
\newcommand{\MainLast}{240}
\newcommand{\MainRowLast}{188}
\newcommand{\IntegerVars}{964}

\title{Shared Pell bounds for modular certificates\\[5pt]
\large Explicit smaller Diophantine equations for ZFC consistency}
\author{Research prepared for Vladimir Reshetnikov's ProveIt project\\
\small AI-assisted derivation, implementation, and independent mathematical checks}
\date{24 October 2026}
\begin{document}
\maketitle
\begin{abstract}
We give a completely specified quartic polynomial in \MainVars\ nonnegative
integer variables, with \MainRows\ quadratic residuals, \MainTerms\ collected
monomials, and coefficient height \MainHeight.  Its no-zero statement is
equivalent to $\Con(\ZFC)$ by an ordinary arithmetic proof composed with the
identified, retained C222--ZF--ZFC refutation calibration.  The construction
encodes finite proof certificates directly; it does not simulate K215 or
instantiate Jones's universal polynomial.

The main change is to extend a single indexed Pell bound to a mixed family
of exact-power and modular-numerator requests.  The enlarged bound is
defined before either kind of request is interpreted.  A separate bootstrap
uses the factorial ceiling equations themselves to justify a smaller digit
radix.  These changes allow one binomial coefficient, one indexed bound,
and one congruence of squared index errors to serve the whole certificate.
The resulting quartic improves all four measured size quantities of the
incoming 261-variable direct quartic.  A smaller-variable, denser companion
and ten-positive-variable norm equations are also supplied.

Every principal residual is printed in the appendix.  The archive includes
the expanded polynomial, executable builders, restoration certificates,
independently written checks, and source proofs.  No consistency assumption
is used to establish the equivalence.  We claim neither a global simplicity
minimum nor a newly completed end-to-end proof-assistant certificate.
\end{abstract}
\clearpage
\tableofcontents
\newpage

\section{Results and the meaning of a simpler equation}

The recommended equation in this packet is
\begin{equation}\label{eq:principal}
 \boxed{P_{\MainVars}(z_0,\ldots,z_{\MainLast})
       =\sum_{j=0}^{\MainRowLast} Q_j(\mathbf z)^2=0,
       \qquad \mathbf z\in\N^{\MainVars}.}
\end{equation}
Here and throughout, $\N=\{0,1,2,\ldots\}$.  The symbols $Q_j$ are the
literal quadratic polynomials in Appendix~\ref{app:residuals}, not
unspecified representations selected by an existence theorem.  The matching
JSON file gives the same ordered rows and every coordinate; the expanded
file gives every coefficient of their collected sum of squares.

\begin{center}
\small
\begin{tabular}{lrrrr}
\toprule
Direct quartic over $\N$ & Variables & Rows & Terms & Height\\
\midrule
Incoming shared-binomial equation &319&254&10,283&86,842,335\\
Incoming shared-bound equation &261&194&31,679&85,566,955\\
\textbf{Principal equation here} &\textbf{\MainVars}&\textbf{\MainRows}&\textbf{\MainTerms}&\textbf{\MainHeight}\\
Dense companion here &\DenseVars&\DenseRows&\DenseTerms&\DenseHeight\\
\bottomrule
\end{tabular}
\end{center}

All four equations have exact degree four.  A ``row'' is one quadratic
residual before the final sum of squares.  ``Terms'' counts distinct
nonzero monomials after collecting that single polynomial.  Height is the
largest absolute value of a collected coefficient.  It is not the
largest coefficient in a residual.  The principal equation saves 20
variables, five residuals, 1,621 monomials, and 379,227 units of coefficient
height relative to the 261-variable equation.  These are reductions of
about $7.66\%$, $2.58\%$, $5.12\%$, and $0.44\%$, respectively.

The dense companion lowers the variable count to \DenseVars, at a substantial
cost in expanded support.  It is included so that minimizing variables
does not silently become minimizing every measure at once.

The comparison is against the complete incoming files at repository commit
\pathfile{f4d8db99e060339f3306a04d507c35bc1d0d68c5}, including all eight ZIPs
then present in \pathfile{docs/incoming}.  The two principal parents are
\emph{Shared Binomial Certificates} and \emph{Shared Pell Bound 261}
\cite{binomial,cap261}.  The repository already has equations with far
fewer variables and very large fixed coefficients, including nine-positive
constructions and a specialized 37-integer quartic.  The results here
advance the explicit direct-calculus, comparatively small-coefficient
branch.  They do not supersede those different tradeoffs.

This distinction is also standard in the literature: degree, number of
unknowns, and the choice between $\N$ and $\Z$ are separate complexity
parameters \cite{bayerdavid,bayeretal}.  None of the present numerical
comparisons is a theorem that no better polynomial exists.

\section{The precise logical statement}

Fix the ordinary finite proof calculus for ZFC used by ProveIt's retained
calibration.  Let $\mathrm{Ref}_{\ZFC}(p)$ be its primitive-recursive
predicate that $p$ codes a finite refutation.  Thus
\[
 \Con(\ZFC)\equiv\neg\exists p\in\N\;\mathrm{Ref}_{\ZFC}(p).
\]
The arithmetic argument of this packet, composed with the retained source
calibration, proves
\begin{equation}\label{eq:logical}
 \ZFC\vdash
 \left(\exists\mathbf z\in\N^{\MainVars}\;P_{\MainVars}(\mathbf z)=0\right)
 \longleftrightarrow
 \left(\exists p\in\N\;\mathrm{Ref}_{\ZFC}(p)\right).
\end{equation}
Taking negations yields exactly the requested consistency equivalence.
The proof is uniform finite arithmetic, rather than an argument restricted
to particular external numerals.  Its ingredients are induction on finite
sequences, elementary Pell theory, integer divisibility, the Chinese
remainder theorem, and explicit finite polynomial substitution.  They
can be carried out on the natural numbers of a model of ZFC.

There are three different levels of evidence, and they should be kept
separate.  First, the literal polynomial and its exact metrics are checked
by finite computation.  Second, the new transformations have ordinary
proofs of soundness and constructive completeness.  Third, the endpoint
from the numerical C222 calculus to ZFC refutations is imported from the
identified repository calibration, whose source texts and pins are retained.
This work does not newly re-prove every transitive calibration dependency
or construct a complete Lean/Rocq proof term for \eqref{eq:logical}.

The imported endpoint is the ordinary argument associated with
\pathfile{KEndpoint.c222_iff_zf_refutation} and
\pathfile{ConstructibleChoice.zf_refutation_iff_zfc_refutation}.
We use that calculus endpoint before any machine-specific endpoint.
Neither the builder nor the definition of the direct quartic reads the
K215 transition table.  The numerical coefficients in
\eqref{eq:principal} are fixed by the supplied finite code and data; their
definition does not branch on whether ZFC is consistent.

\begin{remark}[Unrestricted integer interpretation]
The principal displayed domain is $\N$, not all of $\Z$.  If an equation
quantified over unrestricted integers is required, there is an explicit
quartic conversion without any appeal to a universal polynomial.  Retain
each $z_i$ as an integer variable and add three integer variables
$a_i,b_i,c_i$.  Set
\[
 \widehat P(\mathbf z,\mathbf a,\mathbf b,\mathbf c)
 =P_{\MainVars}(\mathbf z)+
 \sum_{i=0}^{\MainLast}(8z_i+1-a_i^2-b_i^2-c_i^2)^2.
\]
Each new row forces $z_i\ge0$.  Conversely every $8z_i+1$, $z_i\ge0$,
is a sum of three integer squares by the three-squares theorem.  Thus
$\widehat P=0$ has an integer solution exactly when
\eqref{eq:principal} has a natural solution.  This is a transparent
\IntegerVars-variable quartic conversion, not a claimed improvement on the
repository's specialized unrestricted-integer equations.
\end{remark}

\section{The direct proof-certificate source}

The inherited source uses finite multisets of numerical proof rules.  Its
truth label is 0 and its contradiction label is 1.  Formula constructors
are ordinary polynomials; for example the Cantor-style implication
constructor is
\[
 I(a,b)=2(a+b)(a+b+1)+4b.
\]
Membership, universal closure, and equality are the consecutive tags
$I(a,b)+1$, $I(a,b)+2$, and $I(a,b)+3$.  The precise nine axiom families,
modus ponens, generalization, and guarded bundle are retained literally
in the source compiler.  This packet changes the arithmetic certificates
used to validate the bounded source, not the calculus.

Choose a horizon $t\ge12$ and an integer $F$ later forced to be $t!$.
Attach weight $1+jF$ to a positive formula label $j\le t$ and weight 1
to label 0.  For a rule with premises $p_1,p_2$ and conclusion $c$, a
pending-product accumulator satisfies
\begin{equation}\label{eq:flow}
 S'(1+Fc)=S(1+Fp_1)(1+Fp_2).
\end{equation}
The endpoints are $S_0=1+F$ and $S_t=1$.  A bounded sequence of such
relations is encoded using explicit beta remainders.

When $F=t!$, weights attached to distinct positive labels at most $t$
are pairwise coprime.  A prime common to $1+jF$ and $1+kF$ would divide
$j-k$, hence be at most $t$ and divide $F$, contradicting its divisibility
of $1+jF$.  Telescoping \eqref{eq:flow} gives a product identity.  Prime
valuations recover exact multiplicity balance: there is one excess
conclusion of label 1 and zero excess for every other positive label.

This balance entails derivability even if the selected rules contain
cycles.  Let $D$ be the closure of the axioms and truth under the finite
selected rules.  If label 1 were outside $D$, every selected rule whose
conclusion lies outside $D$ would have at least one premise outside $D$.
Counting those premises would contradict the excess of one conclusion
outside $D$.  Thus an accepted flow yields a finite C222 refutation.

Conversely, expand a finite proof DAG into a proof tree, choose $t$ above
all labels, parameters, tree size, and prescribed local polynomial
bounds, and pad with truth rules.  Traversing the pending proof multiset
backwards constructs the positive accumulator products.  After $t$ is
fixed, set $F=t!$, choose a sufficiently large multiple $y$ of $F$, and
use CRT to obtain the beta code $x$.  This ordering matters: the large
accumulator values do not have to be bounded by $t$.

The local source has eight encoded columns and one encoded index.  Its
local polynomial $L$ expresses the rule disjunction and beta equations;
its additional polynomial $J$ bounds all possible rule labels.  Their
literal definitions, finite majorant certificates, and exact source
bindings are included.  Their fixed bounds used below are
\[
 E_0=2^{23},\qquad \ell_0=2^{22}.
\]
These constants concern polynomial growth bounds.  They are not hidden
free parameters of \eqref{eq:principal}.

\section{One indexed Pell bound for exact and modular requests}

For an integer $A\ge2$, write
\begin{align*}
 X_0(A)&=1,&Y_0(A)&=0,\\
 X_{n+1}(A)&=AX_n(A)+(A^2-1)Y_n(A),&
 Y_{n+1}(A)&=X_n(A)+AY_n(A).
\end{align*}
Then $X_n(A)^2-(A^2-1)Y_n(A)^2=1$.  All nonnegative integer solutions
of this Pell equation are obtained uniquely in this way.  This follows
by the usual inverse descent; it is also a retained input to the
parents' explicit indexed-Pell graphs.  The ordinates $Y_n(A)$ strictly
increase, and
\begin{equation}\label{eq:pell-congruence}
 X_n(A)\equiv1,\qquad Y_n(A)\equiv n\pmod{A-1}.
\end{equation}

\subsection{The essential change to the bound}

There are two kinds of requests.  An exact-power request has a base $b$,
an exponent $k$, and an output $m$, intended to mean $m=b^k$.  A modular
request has an upper index $r$, intended to supply the residue of
$(q+1)^r$ for a binomial digit.  A bound including only the exact-power
triples is insufficient: a modular upper index can be arbitrarily larger
than those exponents.

Define $\mathcal A$ to be the set of distinct actual natural-valued
atoms in all exact-power triples, together with every modular upper
index.  The new frame is the literal polynomial equation
\begin{equation}\label{eq:W}
 W=2+\sum_{a\in\mathcal A}a.
\end{equation}
Every independent positive seed index $k_i$ belongs to $\mathcal A$.
If there are $m$ such distinct seeds, then
\[
 0<k_i<W,\qquad m\le\sum_i k_i\le W-2.
\]
These facts hold before any power or binomial has been interpreted.
In the selected construction positivity comes from the raw horizon,
coordinate shifts, and modulus rows.  Indices supplied by an exact Pell
addition formula are recorded as derived indices rather than unchecked
extra seeds.

To make the Pell base sufficiently large, use natural $z,L,\sigma,A$ with
\begin{equation}\label{eq:scale}
 L^2-((W+1)^2-1)\bigl(W(z+1)\bigr)^2=1,
 \qquad A=L+\sigma.
\end{equation}
Pell classification gives $L=X_e(W+1)$ and
$W(z+1)=Y_e(W+1)$ for a positive $e$.  Modulo $W$, the latter is
congruent to $e$, so $W\mid e$ and $e\ge W$.  Therefore
\begin{equation}\label{eq:large-A}
 A\ge L\ge(W+1)^W,\qquad A-1>W^3\quad(W\ge3).
\end{equation}
Conversely, $e=W$ supplies $L,z$, and every larger $A$ remains available
through $\sigma$.  This freedom is essential for the later root congruence.

\subsection{The simultaneous index theorem}

\begin{theorem}[Shared bound and one error congruence]\label{thm:cap}
Let $W\ge3$, $A-1>W^3$, $0<k_i<W$, and $m\le W-2$.  Suppose
$V=Y_W(A)$ is certified once.  The following equations in nonnegative
integers force every principal pair to have exactly its requested index:
\begin{align}
 x_i^2-(A^2-1)y_i^2&=1 &&(1\le i\le m),\label{eq:norms}\\
 V&=W+1+s+\sum_{i=1}^m y_i,\label{eq:cap}\\
 \sum_{i=1}^m(y_i-k_i)^2&=(A-1)u.\label{eq:errors}
\end{align}
Specifically, $x_i=X_{k_i}(A)$ and $y_i=Y_{k_i}(A)$.
Conversely the canonical pairs extend to a solution of these equations.
\end{theorem}
\begin{proof}
Write $y_i=Y_{j_i}(A)$ using Pell classification, initially permitting
$j_i=0$.  Equation~\eqref{eq:cap} gives $y_i<V$, hence $0\le j_i<W$.
Reduce \eqref{eq:errors} modulo $A-1$ using \eqref{eq:pell-congruence}:
\[
 A-1\mid E:=\sum_i(j_i-k_i)^2.
\]
But
\[
 0\le E\le m(W-1)^2\le(W-2)(W-1)^2<W^3<A-1.
\]
Consequently $E=0$, and every $j_i=k_i$.  The potentially enormous
ordinates $y_i$ were not assumed to be small residues; the bound applies
to their decoded indices.

For completeness, $Y_{k_i}(A)\ge k_i$ and their differences are multiples
of $A-1$, so the quotient in \eqref{eq:errors} is a nonnegative integer.
Also $y_i\le Y_{W-1}(A)$ and the recurrence gives
$Y_W(A)>(2A-1)Y_{W-1}(A)$.  Since $m\le W-2$ and $A>W+1$,
$Y_W(A)-\sum_i y_i>W$; hence the slack in \eqref{eq:cap} is natural.
\end{proof}

The theorem is inherited in its original exact-power setting from the
261-variable packet.  The new application is its extension to the mixed
family, using the enlarged \eqref{eq:W}.  Omitting the modular upper
indices would invalidate the size argument and is explicitly detected
by the independent checker.

\subsection{The bound itself is a polynomial certificate}

The notation $V=Y_W(A)$ is proof notation, not an arithmetic function
left in the final equation.  Retain one complete indexed-Pell certificate:
\begin{align*}
 U^2-(A^2-1)V^2&=1,& H^2-(A^2-1)J^2&=1,\\
 S^2-(B_c^2-1)R_c^2&=1,& B_c&=1+4V(d+1),\\
 J&=V^2(e+1),&B_c&=A+Hf,\\
 S&=U+Hg,&R_c&=W+4Vh.
\end{align*}
All letters here are natural; intermediate products are compiled into
quadratic graph rows.  The hypothesis $W\le V$ of the retained indexed
Pell lemma follows from \eqref{eq:cap}.  Its soundness gives
$U=X_W(A)$ and $V=Y_W(A)$, and its constructive completeness supplies
the comparison and congruence witnesses.  The private witnesses of this
one anchor are not inserted into $\mathcal A$: doing so would introduce
an unnecessary circular definition.

\section{A smaller radix justified before factorial semantics}

The shared-binomial parent already identifies the same binomial coefficient
as both a factorial certificate and the reciprocal CRT product.  The new
radix uses this sharing together with the raw ceiling rows.

Let $t=t_0+12$ and take the literal source quantities
\[
 B=b_0+2,\quad R=2+t+x+y+F,\quad K=R,\quad
 N=B^K,\quad V_0=(F^2+1)^t.
\]
Let $S_0$ be the sum of the nine positive CRT representatives.  Define
\begin{equation}\label{eq:radix}
 q=N(S_0+1)^K.
\end{equation}
The displayed powers are realized by the simultaneous polynomial
power graph; they are not variable exponents in the final polynomial.
Compared with the earlier self-sized radix, the term $V_0$ has disappeared
from its base.  The aggregate upper bound on the nine representatives is
also unnecessary: their actual sum sets the radix after they are chosen.

The retained ceiling rows are
\begin{equation}\label{eq:ceiling}
 V_0+B=BF+r+1,\qquad B=r+s+1.
\end{equation}
Thus $V_0+s=BF$ and $0\le s<B$.  Once the exact powers are decoded,
$V_0\ge1$, so $F\ge1$.  Since $t\ge12$,
\begin{equation}\label{eq:bootstrap}
 F^2<V_0\le BF\quad\Longrightarrow\quad F<B,
 \qquad V_0<B^2\le B^K=N\le q.
\end{equation}
This uses neither $F=t!$ nor the interpretation of $B$ as a binomial.
The raw inequality $V_0<B^2$ already occurs in the incoming unified-binomial
packet \cite{unified}; its application to the new self-sized radix and
mixed-index certificate is the combination used here.

\begin{lemma}[Uniform lower binomial digits]\label{lem:binomial-bound}
For integers $r\ge1$, $t\ge1$, and $0\le j\le t$,
\[
 \binom{r+t}{j}\le(r+1)^t.
\]
\end{lemma}
\begin{proof}
Induct on $t$.  The case $t=1$ is immediate.  If $j=t$, use
$\binom{r+t}{t}=\prod_{i=1}^t(1+r/i)\le(r+1)^t$.
For $1\le j<t$, Pascal's identity and induction bound its two terms by
$2(r+1)^{t-1}\le(r+1)^t$.  The case $j=0$ is immediate.
\end{proof}
Apply the lemma with $r=F^2$.  All digits through place $t$ in
$(q+1)^{F^2+t}$ are at most $V_0<q$.  The single exact-binomial digit
certificate therefore forces
\begin{equation}\label{eq:shared-B}
 B=\binom{F^2+t}{t}.
\end{equation}
No later CRT range estimate is needed to recognize this first binomial.

\begin{lemma}[Squared-radius factorial fixed point]\label{lem:factorial}
For $t\ge6$ and a positive integer $F$,
\[
 F=\ceil{\frac{(F^2+1)^t}{\binom{F^2+t}{t}}}
 \quad\Longleftrightarrow\quad F=t!.
\]
\end{lemma}
\begin{proof}
Put $H_t(F)=t!\prod_{i=1}^t(F^2+1)/(F^2+i)$.  Each factor is at most
one, so the ceiling relation gives $H_t(F)\le F\le t!$.
If $F<t$, monotonicity in $F^2$ gives
$H_t(F)\ge2^t/(t+1)>t>F$, impossible for $t\ge6$.
For $F\ge t$, the product inequality gives
\[
 H_t(F)>t!\left(1-\frac{t(t-1)}{2F^2}\right)>\frac{t!}{2}.
\]
If $F<t^2$, then $t!/2>t^2$ gives another contradiction.  Thus
$F\ge t^2$, and
$H_t(F)>t!(1-1/(2F))$.  If $t!\ge F+1$, the last lower bound is at
least $(F+1)(1-1/(2F))\ge F$, again impossible.  Hence $F=t!$.
Conversely at $F=t!$ the same estimates give
$F-1/2<H_t(F)\le F$, proving the ceiling relation.
\end{proof}
This factorial lemma is prior work in the repository and the
shared-binomial parent; it is reproduced to make the dependency order
transparent.  Equations~\eqref{eq:ceiling} and \eqref{eq:shared-B} now
give $F=t!$.

\subsection{All later digits and the smaller amplification}

Every ordinary upper representative $a$ satisfies $a\le S_0$.  Every
requested lower index $k$ is less than $R=K$.  Thus all its required
binomial coefficients are bounded by
\[
 a^k\le(S_0+1)^k<(S_0+1)^K\le q.
\]
For the derived label polynomial, the retained literal positive-DAG
majorant is $J\le(S_0+1)^{\ell_0}$.  At $t=12$,
\[
 \ell_0(t+1)=2^{22}\cdot13<12!=479,001,600.
\]
The inequality persists under increasing $t$, so
$\ell_0(t+1)<t!=F<R=K$.  Hence its digits through place $t+1$ are also
strictly below $q$.

The local evaluation bound is less than $R^{E_0}$.  For $R\ge12!>2^{28}$,
\[
 R^{2^{23}}<2^R.
\]
One elementary verification takes $2^m\le R<2^{m+1}$, $m\ge28$.
The inequality $(m+1)2^{23}<2^m$ holds at $m=28$ and persists by induction.
It follows that $R^{2^{23}}<2^{(m+1)2^{23}}<2^R$.
Consequently amplification $K=R$ is sufficient; the former factor of two
is unnecessary on this horizon.  These estimates occur after the factorial
has been recovered, while \eqref{eq:bootstrap} occurs before it.

\section{Modular numerators, exact powers, and CRT decoding}

\subsection{Recovering all exact powers}

For a retained request $m=b^k$, introduce the nonnegative modulus
$T_b=2Ab-b^2-1$ and impose
\[
 X_k(A)+bY_k(A)=AY_k(A)+m+T_bh,\qquad h\ge0.
\]
The raw modulus row excludes $b=0$.  The mixed frame places $b,k,m<W$,
and \eqref{eq:large-A} makes $T_b$ larger than both $m$ and $b^k$.
The Pell recurrence modulo $T_b$ gives
\[
 X_k(A)+(b-A)Y_k(A)\equiv b^k\pmod{T_b}.
\]
By Theorem~\ref{thm:cap}, this is the actual requested $k$, not an
unspecified congruent index.  The strict interval then forces $m=b^k$.
The positive expression on the left also supplies the required natural
quotient in the constructive converse.  Thus the exact powers can be
decoded before the first binomial is interpreted.

\subsection{One freely chosen root for every numerator}

Let $b=q+1$.  The retained common modulus $M$ is a product of the digit
denominator powers and $q$, large enough that every required $qd$ and
$Nd$ divides $M$.  Its prime divisors divide $q$.  Add
\begin{equation}\label{eq:root}
 2A(q+1)=(q+1)^2+1+Mh_0,\qquad h_0\ge0.
\end{equation}
Then $s=q+1-A$ satisfies $s^2\equiv A^2-1\pmod M$.  The identity
$(A+s)^r=X_r(A)+sY_r(A)$ in the quadratic algebra yields
\begin{equation}\label{eq:residue}
 \mathcal R_r:=X_r(A)+(q+1-A)Y_r(A)
       \equiv(q+1)^r\pmod M.
\end{equation}
Moreover $\mathcal R_r>0$ for $r>0$, because
$X_r(A)>(A-1)Y_r(A)$ and $q>0$.

For the exact binomial at upper index $r=F^2+t$ and denominator $d=q^t$,
use
\[
 \mathcal R_r=r_0+dB+qd h,\qquad
 d=r_0+s_0+1,\qquad q=B+o+1.
\]
For a range request $N\mid\binom{r}{k}$, use
\[
 \mathcal R_r=r_0+Nd h,\qquad d=q^k=r_0+s_0+1.
\]
The carry bounds of the preceding section and the binomial theorem
identify these as the desired exact digit and divisibility conditions.
The upper index of every such request is one of the newly included atoms
in \eqref{eq:W}, so its Pell pair has already been certified by the same
single anchor.

Completeness of \eqref{eq:root} is not automatic from its notation.
Since $\gcd(q+1,M)=1$, the gcd of $2(q+1)$ and $M$ is at most 2.
If it is 2, $q+1$ is odd and $(q+1)^2+1$ is even.  Thus the linear
congruence for $A$ is soluble.  Its solution class contains arbitrarily
large positive integers.  After all semantic data, $W$, and $L$ are fixed,
choose such an $A\ge L$ large enough that $h_0\ge0$.  All Pell and cap
witnesses are then constructed at this chosen base.  The positive residues
are congruent to the least prescribed residues modulo their moduli, so
the remaining quotients are nonnegative.

\subsection{The reciprocal product and bounded source}

Now $F=t!$ and $B=\binom{F^2+t}{t}$.  For $1\le i\le t$, define
$d_i=1+F^2/i$.  They are integers, pairwise coprime, their prime factors
exceed $t$, and
\[
 \prod_{i=1}^t d_i=B.
\]
The retained CRT source imposes, for its encoded index $C$ and eight
encoded columns $A_j$,
\[
 B\mid C+F^2+1,\qquad N\mid\binom C t,
 \qquad N\mid\binom{A_j}{b_j+1},
\]
\[
 N\mid L(C,\mathbf A),\qquad
 N\mid\binom{J(\mathbf A)}{t+1}.
\]
The one-way binomial range lemma states that
$p^e\mid\binom ak$ implies $a\bmod p^e<k$.  For $a\ge k$, its proof
uses the usual floor-sum valuation formula: if $[a-k+1,a]$ contains no
multiple of $p^e$, the valuation is at most $e-1$.  The case $a<k$ is
immediate.  The false converse is never used.

Apply this at a prime factor of each $d_i$, with precision multiplied by
$K$.  It recovers the index $i-1$ and the bounded local witnesses.  Since
$2^K>R^{E_0}$, divisibility of the nonnegative local evaluation forces its
actual value to be zero.  The derived range yields the label bounds.
This reconstructs every bounded source row and hence the finite refutation
described in Section~3.

In the converse direction, use CRT at twice the valuation precision.
For each prime $p$ of $d_i$, the numerator of the relevant falling
factorial has a factor divisible by $p^{2K v_p(d_i)}$.  Division by $k!$
loses at most $k-1<K$ in valuation, leaving the required precision.
Choose positive representatives and only then form their sum $S_0$ and
the radix \eqref{eq:radix}.  This removes the old need to fit those
representatives into an independently chosen aggregate bound.

\section{Literal polynomial compilation and verification}

Every formula above is translated into a finite quadratic graph using
addition and multiplication nodes, norm rows, remainder equations, and
the fixed local calculus polynomial.  No factorial, exponentiation,
binomial coefficient, Pell function, bounded quantifier, or divisibility
symbol remains in the final list of residuals.

The principal source contains 777 semantic coordinates and 725 rows.
An independently checked affine export restores 428 eliminated coordinates
by polynomials known to be nonnegative on the natural domain.  This gives
349 coordinates and 297 quadratic residuals.  A further exact certificate
eliminates 108 coordinates, leaving \MainVars\ coordinates and \MainRows\
residuals.  Each elimination has a displayed restoration formula and a
literal polynomial check.  Preserving natural-valued restorations is
essential: an arbitrary integer-valued substitution would not preserve
the stipulated domain.

The compiler's support threshold is an implementation choice controlling
which quadratic substitutions are retained.  The selected threshold 32
gives the principal equation.  Larger thresholds give the denser
\DenseVars-variable companion.  The algorithm is not claimed to find an
optimal elimination order or a globally minimal graph.

All remaining residuals have degree at most two.  At least one has a
nonzero quadratic part, so the sum of their squares has exact degree four:
over the reals a sum of squares of homogeneous quadratic forms cannot
vanish identically unless every form is zero.  The fully collected
expansion independently confirms the degree and gives the exact term
and height counts in Section~1.

The checks supplied with the packet cover the literal source rows, mixed
frame atom set, anchor certificate, positive independent indices, pooled
error congruence, ordinary polynomial expansion, and every deleted-variable
restoration.  Rebuilding the selected files in a fresh directory must give
the same mathematical artifacts.  Random numerical evaluations are used
only as supplementary transcription checks where indicated; they are not
the proof of an unbounded arithmetic claim or of the consistency endpoint.

\section{A short ten-positive-variable companion}\label{sec:positive}

The direct quartic is the principal result.  A separate improvement on the
compact native branch uses a coordinate condition already forced at every
zero.  This branch retains the native/K215 calibration and its four large
fixed coefficients, unlike the direct quartic.

Let $(\beta,K_0,M_C,M_F)$ be the complete original compact44 tuple supplied
in \pathfile{positive/coefficients}.  These are fixed integers, not
unknowns.  Their raw binary sizes are 590,069, 592,884, 590,069, and 590,070
bits respectively; $\beta=2^{590069}-1$.  The entire tuple is retained,
with $M_F$ already shifted exactly once.  Its payloads and finite
reconstruction data are supplied together.

The ten unknowns are
\[
 J,F,\alpha,i,T,w,s,n,\eta,\zeta\in\Pos.
\]
Define the following polynomial abbreviations:
\begin{align*}
 q&=\beta J+1,&g&=q-1,&X&=wq,&Y&=q^3s,\\
 \mu&=q^2+1,&k&=\eta+\zeta,&c&=kY+\eta,\\
 A&=\mu Y(X+1)+1,&\Delta&=A^2-1,&H&=4A-5,\\
 V_1&=X\mu^2Y^2,&D&=V_1(V_1+1)k^2+1,\\
 C&=q-F-\alpha,&B&=(K_0+w)C-F,\\
 G&=(q-1)(q-F)+\alpha,&
 R&=G(q^2-1)+(M_C+qM_F)J,\\
 e&=k-R-1,&E&=ic^2,\\
 F_s&=1+\Delta DE^2,&Q&=\Delta^2DE^2,&\delta&=Q-1,\\
 b&=c(T+\alpha)-R,&V&=F_sb-c,&D_*&=QV^2-1,\\
 K_+&=X(3c-2X)+2,&M&=g\mu^2H,&a&=DF_s.
\end{align*}
For the shortest version put
\begin{equation}\label{eq:positive-Z}
 Z=g\mu^2\bigl[H(n+V-bB)+K_+\bigr]
       -H\bigl(\mu^2B+ge\bigr).
\end{equation}
The equation is the explicit polynomial
\begin{equation}\label{eq:positive-N}
 \boxed{\mathcal N=
 \bigl(\delta Z^2-M^2D_*-\delta M^2a\bigr)^2
       -4\delta M^4aD_*=0.}
\end{equation}
Every exponent in these definitions is a fixed nonnegative integer.
There are no radicals or divisions in the equation.  The fully paid
division-free schedule takes \textbf{86 binary operations}: 51
multiplications, 18 additions, and 17 subtractions, including constant
multiplications and squaring.  Its exact total degree is \textbf{466}.

\subsection{Quotient chart and exact cancellation}

The pinned uncentered gap-fusion parent takes 89 operations and has degree
288 \cite{gapfusion}.  At every positive zero its proof yields $q^3\mid Y$.
Substitute $Y=q^3s$ into that polynomial.  Both its carrier numerator and
denominator acquire a factor $q^3$, so its quartic norm acquires $q^{12}$.
Canceling that explicit nonzero factor yields a polynomial of degree 466
with 87 operations.  This version is in exact bijection with the parent's
positive zero set: it retains $n+s$ in place of $n$ in
\eqref{eq:positive-Z}.

The shortest 86-operation version absorbs the integer term $s$ into the
positive padding coordinate.  The exact universal polynomial identity is
\begin{equation}\label{eq:quotient-identity}
 N_{\rm parent}(Y=q^3s,\ n_{\rm old}=n-s)
       =q^{12}\mathcal N.
\end{equation}
The identity alone is not enough for soundness because $n-s$ could be
nonpositive.  The following additional argument handles that point.

\subsection{Why absorbing the padding is sound}

All pre-transport steps of the parent proof are independent of the sign
of its padding variable.  At a positive zero of \eqref{eq:positive-N}
they recover positive integer roots $\tau,f,y$ with
\[
 D=\tau^2,\quad F_s=f^2,\quad
 QV^2-(Q-1)y^2=1,
\]
and give the early estimates
\[
 b>8Q,\qquad \tau f<b/2,\qquad 0<y-V<b/4,
 \qquad e/\mu^2<b/4,\qquad K_+>0.
\]
The fifth-index gap underlying $b>8Q$ is prior work of the parent, not a
new lemma here.  The two-root norm gives, for signs $\varepsilon,\nu$,
\[
 \frac ZM=\varepsilon\tau f+\nu y,
 \qquad
 \frac ZM=n+V+\frac{K_+}{H}
       -\frac{(1+gb)B}{g}-\frac e{\mu^2}.
\]
Therefore
\[
 n=\varepsilon\tau f+(\nu y-V)-\frac{K_+}{H}
       +\frac{(1+gb)B}{g}+\frac e{\mu^2}
   <b+\frac{(1+gb)B}{g}.
\]
If $B\le-1$, the right side is at most $-1/g<0$, contradicting
$n>0$.  Thus the same transport sign $B\ge0$ is forced.

The subsequent rank, index, parity, projection, and source-reconstruction
arguments use the integral carrier and this sign.  The condition $q^3\mid Y$
is now built into the coordinates; the remaining denominator is
$g\mu^2H$, strengthening the relevant $H$-divisibility statement.  The
whole proof recovers the same native source.  Conversely a canonical
parent witness has $Y=q^3s$ with $s>0$, and setting
$n=n_{\rm old}+s$ is positive and satisfies \eqref{eq:quotient-identity}.
This is an existence-preserving reparameterization.  We do not assert
that its unshifted inverse sends every zero to a positive old padding
coordinate.

\subsection{Exact degree and the tradeoff}

All ten unknowns have degree one, while the coefficient ports have degree
zero.  In the quotient chart the relevant degrees are
\[
 \begin{array}{c|rrrrrrrrrrrr}
 \text{register}&Y&A&\Delta&D&E&F_s&Q&b&V&M&a&Z_0\\\hline
 \text{degree}&4&8&16&30&11&68&84&6&74&13&98&\le21,
 \end{array}
\]
where $Z=MV+Z_0$.  Expand the norm after this substitution.  Its unique
highest-degree contribution is
\[
 -4\delta^2M^4V^2a,
\]
with degree $2\cdot84+4\cdot13+2\cdot74+98=466$.  More explicitly its
nonzero leading homogeneous form is
\[
 -2^{10}(\beta J)^{326}i^{10}w^{30}s^{64}
       (\eta+\zeta)^{34}(T+\alpha)^2.
\]
The leading form of $b$ is
$(\eta+\zeta)Y_0(T+\alpha)$, not the old leading form $-R_0$; this change
must be included when checking the new degree.

Adding the sharp product center $ADE$ gives companion schedules with
degree 450 and 89 operations (absorbed padding) or 90 operations (exact
parent-positive bijection).  The factor $\Gamma^2-a=D^2E^2-D$ lowers
the leading contribution by 16.  These are diagnostic controls: the
89-operation, degree-450 version is dominated by the already existing
89-operation, degree-288 equation.  The 86-operation equation improves
arithmetic length at the same ten-positive-variable count.

\IfFileExists{centering.tex}{\section{A ten-positive-variable equation of degree 246}\label{sec:center246}

A second native companion improves degree rather than operation count.
It has ten positive unknowns, exact degree \textbf{246}, and a completely
paid schedule of \textbf{121 operations} (73 multiplications, 25 additions,
and 23 subtractions).  Its parent sharp-center equation has degree 254
and 97 operations \cite{gapfusion}.  This is a degree/cost tradeoff.
The following construction uses the same complete compact44 tuple as
Section~\ref{sec:positive}, but restores the independent coordinate $Y$:
\emph{the substitution $Y=q^3s$ is not made here}.

\subsection{The full polynomial and its two correction terms}

Let $J,F,\alpha,i,T,x,Y,n,\eta,\zeta\in\Pos$.  Define
\begin{align*}
 q&=\beta J+1,&X&=q+x,&g&=q(q-1),\\
 C_d&=q-F-\alpha,&B&=(qK_0+X)C_d-qF,\\
 \mu&=q^2+1,&t&=\mu Y,&L&=X+1,\\
 A&=tL+1,&\Delta&=A^2-1,&H&=4A-5,\\
 k&=\eta+\zeta,&c&=kY+\eta,&v&=Xt^2,\\
 D&=1+v(v+1)k^2,&E&=ic^2,\\
 F_s&=1+\Delta DE^2,&Q&=\Delta^2DE^2,&\delta&=Q-1.
\end{align*}
The remaining unchanged kernel and carrier data are
\begin{align*}
 R&=[(q-1)(q-F)+\alpha](q^2-1)+(M_C+qM_F)J,\\
 e&=k-R-1,&b&=c(T+\alpha)-R,&V&=F_sb-c,\\
 D_*&=QV^2-1,&P&=q^4(q-1),&M&=P\mu^2H,\\
 K_+&=X(3c-2X)+2,&W&=1+gb.
\end{align*}
The new polynomial abbreviations are
\begin{align}
 Q_1&=k^2\bigl[t^2LX^2(tL-1)+t(L^3-2L+1)-L^2+2\bigr],\label{eq:center-Q1}\\
 Q_3&=k^2(tX^2-3L+6),&\lambda&=8L^3,\label{eq:center-Q3}\\
 \Gamma&=E(8L^3AD-4Q_1-Q_3),&a&=\lambda^2DF_s.\label{eq:center-Gamma}
\end{align}
Set
\begin{equation}\label{eq:center-Z}
 \begin{aligned}
 Z={}&P\mu^2\bigl[H(n+V+\Gamma-bB)+K_+\bigr]\\
     &+H\bigl[\mu^2(gY-q^3B-(q-1)X)-Pe\bigr].
 \end{aligned}
\end{equation}
The complete equation is
\begin{equation}\label{eq:center246}
 \boxed{\mathcal N_{246}
 =\bigl(\delta Z^2-M^2D_*-\delta M^2a\bigr)^2
       -4\delta M^4aD_*=0.}
\end{equation}
All abbreviations are integer polynomials in the ten listed unknowns.
The corresponding literal graph and coefficient payloads are in
\pathfile{center246/}.  No additional quantified coordinate is hidden
in~\eqref{eq:center246}.

The construction comes from two exact polynomial divisions.  With
$\Omega=L^3D$ and $X=L-1$, expansion gives
\begin{align}
 \Omega&=AQ_1+r_1=A^3Q_3+r_3,\label{eq:center-divisions}\\
 r_1&=L^3+k^2(L^2-t-2),\nonumber\\
 r_3&=L^3+k^2\bigl[L^2t^2(L^2+5L-3t-12)
             +Lt(8L-3t-16)+3L-t-6\bigr].\nonumber
\end{align}
Thus $8\Omega A-4Q_1-Q_3$ retains the polynomial portions of the
first three terms of
$8\Omega\sqrt{A^2-1}=8\Omega A-4\Omega/A-\Omega/A^3-\cdots$.
The series is a discovery guide only: the construction and proof use
the exact identities~\eqref{eq:center-divisions}.
For a shorter evaluation schedule, put $p=A-1$, $u=tX^2$ and use
\begin{equation}\label{eq:center-paid}
 4Q_1+Q_3=k^2\bigl[u\{4p(p-1)+4L+9\}
                    +4(tX-L^2)-3L+14\bigr].
\end{equation}

\subsection{Positive bounds, root recovery, and transport}

\begin{lemma}\label{lem:center-bounds}
On the whole positive input domain, $0<\Gamma<Q$.
At every positive zero of~\eqref{eq:center246}, the recovered Pell roots
satisfy $\lambda\tau f+X/q^4<2Q<b/4$.
\end{lemma}
\begin{proof}
Here $L\ge33$ and $t\ge1025$.  Both $Q_1,Q_3$ are positive:
the last part of $Q_1/k^2$ is at least
$L^3-L^2-2L+3>0$, while
$Q_3/k^2\ge(L-2)(L-3)+1>0$.
The first term of $Q_1/k^2$ is at least $24t^3>t+2$, so
$r_1>-k^2(t+2)>-Q_1$ and~\eqref{eq:center-divisions} gives
$\Omega>(A-1)Q_1>Q_1$.  Also $Q_3<k^2tX^2<D\le\Omega$.
It follows that
\[
 0<\Gamma<\lambda ADE,\qquad
 \lambda A=8L^3(tL+1)\le16tL^4<t^4L^4<\Delta^2.
\]
Since $E\ge1$, this proves $\Gamma<Q$.

For completeness, root recovery does not require the desired transport
sign.  Since $F_s>c+1$ and $b$ is integral, $|V|>1$ and $D_*/\delta>0$.
The positive radicands $a$ and $D_*/\delta$ are unequal: equality
would give $\delta\lambda^2DF_s=QV^2-1$, impossible modulo $D>1$.
Writing $z=Z/M$, the norm therefore has $z\ne0$ and gives the rational
square root $(z^2+a-D_*/\delta)/(2z)$ of $a$.
As $\lambda>0$ and $\gcd(D,F_s)=1$, it follows that
$D=\tau^2$ and $F_s=f^2$ for positive integers $\tau,f$.
With $I=\tau E$, the unchanged adjacent-Pell descent gives an integer
$y>0$ and
\[
 f^2-\Delta I^2=1,\qquad Q=(\Delta I)^2,\qquad
 QV^2-(Q-1)y^2=1,
 \qquad z=\varepsilon\lambda\tau f+\nu y.
\]
Here $\varepsilon,\nu\in\{1,-1\}$.
The parent's pre-transport folding, first-index, and fifth-index
arguments now give $b>8Q$, $K_+>0$,
$0<y-V<b/4$, and $e/\mu^2<b/4$ \cite{gapfusion}.
None uses the new integral center or radical weight.
Finally $f<AI$ implies $\lambda\tau f<\lambda ADE<Q$;
also $X/q^4<Q$.  Their sum is less than $2Q<b/4$.
\end{proof}

The exact normalized carrier is
\begin{equation}\label{eq:center-carrier}
 \frac ZM=n+V+\Gamma+\frac{K_+}{H}+\frac Y{q^3}
          -\frac{WB}{g}-\frac X{q^4}-\frac e{\mu^2}.
\end{equation}
Combining it with the recovered signed roots and
Lemma~\ref{lem:center-bounds} gives
\[
 n<\frac{3b}{4}+\frac{WB}{g}.
\]
If $B\le-1$, then $WB/g\le-b-1/g$, contradicting $n>0$.
Thus $B\ge0$ and $C_d>0$.  The remaining ordered rank, parity,
coprimality, whole-product projection, and unscaled source reconstruction
are those of the pinned gap-fusion proof.  The center and weighted root
are integral; the fractional terms in~\eqref{eq:center-carrier} and
$W\equiv1\pmod g$ are unchanged.  Hence the same source61/native/controller
calibration recovers a finite ZFC refutation.  This is the inherited
native branch, not the direct C222 route of the principal quartic.

Conversely take the parent's fresh canonical kernel witnesses from a
refutation.  They supply integral $t_0=B/g\ge0$, $\rho=e/\mu^2\ge1$,
$s_Y=Y/q^3>0$, $\kappa=K_+/H>0$, $X/q^4>0$, and
$u=y-V>2Q$, with $s_Y,\kappa<c$.  Define the remaining coordinate by
\begin{equation}\label{eq:center-complete}
 n=Wt_0+\rho+\lambda\tau f-\Gamma+u-s_Y-\kappa+X/q^4.
\end{equation}
It is an integer and is positive because
$n>2Q-Q-2c>0$; here $Q>c^4>2c$ follows from $E=ic^2$ and $c\ge3$.
The carrier becomes $\lambda\tau f+y$, so the norm vanishes.
This proves existence equivalence by a fresh positive construction,
without claiming a positive pointwise inverse for every parent zero.

\subsection{Exact degree and a noncancelling coefficient}

Write $C_*:=\Gamma^2-a$.  Assign weights $1,3,1,5$ to the independent
polynomial variables $L,t,k,E$.  Then
$\degr A=4$, $\degr D=16$, $\degr\Omega=19$,
$\degr Q_1=15$, $\degr Q_3=7$, $\degr r_1=5$, and $\degr r_3=13$.
An exact identity derived from~\eqref{eq:center-divisions} is
\begin{equation}\label{eq:center-defect}
 \begin{aligned}
 A^2(C_*+64L^6D)=E^2\bigl[&64A^2\Omega r_1
       +16\Omega(r_3-2r_1)+16r_1^2\\
       &+8A^2Q_1Q_3+A^2Q_3^2\bigr].
 \end{aligned}
\end{equation}
The two degree-42 terms on the right have combined leading part
$-112L^7t^7k^4E^2$.  All others have smaller degree.
Since $A_0=Lt$, this proves
\begin{equation}\label{eq:center-defect-leading}
 \degr C_*=34,\qquad (C_*)_0=-112L_0^5t_0^5k_0^4E_0^2.
\end{equation}
The supplied certificate independently lists and checks all 52
coefficients of $C_*$.

In the actual ten inputs, the unchanged kernel has
\[
 \degr\delta=42,\quad\degr V=38,\quad\degr M=13,
 \quad\degr a=56,\quad\degr\Gamma=28.
\]
Write $Z=M(V+\Gamma)+Z_0$, where $\degr Z_0\le19$ and
$(Z_0)_0=-M_0b_0B_0$.  Put
\[
 \Theta=\delta M^2C_*-M^2(V^2-1)
       +2\delta M(V+\Gamma)Z_0+\delta Z_0^2.
\]
Direct expansion gives
\begin{equation}\label{eq:center-norm-expansion}
 \begin{aligned}
 \mathcal N_{246}={}&4\delta^2M^4V^2C_*
       +4\delta M^2V\Gamma\Theta+\Theta^2\\
       &-4\delta M^4aV^2+4\delta M^4a.
 \end{aligned}
\end{equation}
Here $\degr\Theta=112$, so the five displayed terms have degrees at most
$246,246,224,226,150$.  The first two terms both contribute at degree 246;
their possible cancellation must be checked.

Let $q_0=\beta J$, $X_0=q_0+x$, $t_0=q_0^2Y$, $k_0=\eta+\zeta$,
$E_0=ik_0^2Y^2$, $b_0=-q_0^3(q_0-F)$, and
$B_0=(q_0K_0+X_0)(q_0-F-\alpha)-q_0F$.
The complete leading form is
\[
 4\delta_0^2M_0^4V_0^2
 \bigl((C_*)_0-2\Gamma_0b_0B_0\bigr),
 \qquad\Gamma_0=8X_0^6t_0^5k_0^2E_0.
\]
Its first summand has degree ten in $i$ and equals
\[
 -7\cdot2^{14}q_0^{116}X_0^{29}i^{10}Y^{57}
                     k_0^{32}(q_0-F)^2;
\]
the second has degree nine in $i$.  Consequently
\begin{equation}\label{eq:center-coefficient}
 \boxed{[J^{116}x^{29}i^{10}Y^{57}\eta^{32}F^2]\,
          \mathcal N_{246}=-7\cdot2^{14}\beta^{116}\ne0.}
\end{equation}
This proves exact degree 246 for the actual fixed tuple, rather than
only for generic ports.

The companion checker passes 151 checks, including exact local polynomial
identities, complete coefficient and source pins, circuit liveness,
independent formula evaluations, and degree/leading-form transcription
checks over three finite fields.  Polynomial-part radical centering
already appears elsewhere in ProveIt, for example the unrestricted-integer
construction in \pathfile{87-direct-centered-round12/norm/PROOF.md}.
The new result is this concrete application and degree/cost point relative
to the frozen snapshot.  It inherits the parent's ordinary calibration;
no new end-to-end proof-assistant certificate or global minimality claim
is made.
}{}

\section{What is new, what is inherited, and what remains open}

The contribution of this round is a proved and fully materialized
combination, together with a new use of a raw bootstrap inequality.  The single anchor
must include modular indices, the self-sized radix is justified by
$V_0\le BF$ before factorial semantics, and the complete contracted
quartic is counted and supplied.  Merely placing two prior algorithms
next to one another would not establish any of those points.

The factorial fixed-point lemma, common modular Pell residues, binomial
range method, reciprocal CRT source, one-anchor error-square theorem,
and natural-variable contraction each have explicit predecessors in
the retained sources.  The relation-combining and Pell tradition is
classical.  We make no claim to have invented those general techniques.
The novelty audit is bounded to the inspected repository snapshot,
incoming archives, and cited literature; it is not a worldwide
priority search or a minimum-size proof.

The native companion establishes a different saving: a forced quotient
condition can be built into the coordinates, a homogeneous clearing
factor can be removed, and an integer term can be absorbed into a padding
coordinate after re-proving its sign argument.  The distinction between
an exact positive-zero bijection and fresh positive completeness is
essential in that last step.

There are still concrete opportunities.  A better contraction order may
reduce the direct quartic without changing the number theory.  A smaller
direct proof calculus may simplify its local source, provided its full
refutation equivalence is proved.  A common indexed bound tailored to the
actual modular indices might reduce the remaining anchor overhead.
None of these possibilities is included in the proved counts above.

The packet establishes an equivalence, not either side of it.  In
particular it neither finds a refutation nor proves that all tuples fail
to solve the equation.  A newly checked, end-to-end formal proof remains
a separate task.  The delivered finite objects and explicit proof
boundaries are intended to make that task and further optimization
concrete.

\appendix
\section{Every residual of the principal quartic}\label{app:residuals}

The following is the complete definition of the polynomials $Q_j$ in
\eqref{eq:principal}.  Indices are zero-based.  All $z_i$ independently
range over $\N$.  Line breaks within one $Q_j$ are typographical only.
The appendix is generated directly from the authoritative equation file
and checked against it by the release verifier.
\begingroup\footnotesize
\setlength{\tabcolsep}{3pt}
\begin{longtable}{@{}r p{0.88\linewidth}@{}}
\toprule Residual & Integer polynomial (continued across rows when necessary)\\\midrule\endhead
$Q_{0}=$ & $-z_{2}^{2}\,+z_{3}$\\[5pt]
$Q_{1}=$ & $1\,-2z_{2}\,-z_{2}z_{4}\,+z_{4}\,+z_{5}\,-z_{6}$\\[5pt]
$Q_{2}=$ & $1\,+z_{4}\,-z_{6}\,-z_{7}$\\[5pt]
$Q_{3}=$ & $-13\,-z_{0}\,-z_{0}z_{2}\,-z_{0}z_{9}\,-13z_{2}\,-13z_{9}$\\
 & $+z_{11}$\\[5pt]
$Q_{4}=$ & $-1\,+z_{1}\,-z_{2}\,-z_{2}z_{8}\,-2z_{8}\,-z_{8}z_{9}$\\[5pt]
$Q_{5}=$ & $-1\,+z_{1}\,-z_{10}\,-z_{10}z_{11}$\\[5pt]
$Q_{6}=$ & $2\,+z_{3}\,-z_{4}z_{22}\,+z_{20}\,-2z_{22}$\\[5pt]
$Q_{7}=$ & $-6\,-5z_{12}\,-z_{12}^{2}\,-2z_{12}z_{13}\,-5z_{13}\,-z_{13}^{2}$\\
 & $+z_{23}$\\[5pt]
$Q_{8}=$ & $-72\,-68z_{13}\,-16z_{13}^{2}\,-16z_{13}z_{23}\,-34z_{23}\,-4z_{23}^{2}$\\
 & $+z_{24}$\\[5pt]
$Q_{9}=$ & $-30\,-11z_{12}\,-z_{12}^{2}\,-4z_{12}z_{24}\,-22z_{24}\,-4z_{24}^{2}$\\
 & $+z_{25}$\\[5pt]
$Q_{10}=$ & $-380\,-312z_{24}\,-64z_{24}^{2}\,-32z_{24}z_{25}\,-78z_{25}\,-4z_{25}^{2}$\\
 & $+z_{26}$\\[5pt]
$Q_{11}=$ & $-6\,-5z_{13}\,-z_{13}^{2}\,-2z_{13}z_{14}\,-5z_{14}\,-z_{14}^{2}$\\
 & $+z_{27}$\\[5pt]
$Q_{12}=$ & $-6\,-5z_{12}\,-z_{12}^{2}\,-2z_{12}z_{14}\,-5z_{14}\,-z_{14}^{2}$\\
 & $+z_{28}$\\[5pt]
$Q_{13}=$ & $-414\,-22z_{12}\,-2z_{12}^{2}\,-16z_{12}z_{14}\,-8z_{12}z_{27}\,-200z_{13}$\\
 & $-32z_{13}^{2}\,-64z_{13}z_{14}\,-32z_{13}z_{23}\,-32z_{13}z_{28}\,-320z_{14}\,-64z_{14}^{2}$\\
 & $-32z_{14}z_{23}\,-32z_{14}z_{27}\,-32z_{14}z_{28}\,-100z_{23}\,-8z_{23}^{2}\,-16z_{23}z_{28}$\\
 & $-52z_{27}\,-8z_{27}^{2}\,-108z_{28}\,-8z_{28}^{2}\,+z_{29}$\\[5pt]
$Q_{14}=$ & $-1344\,-800z_{13}\,-128z_{13}^{2}\,-256z_{13}z_{14}\,-128z_{13}z_{23}\,-128z_{13}z_{28}$\\
 & $-864z_{14}\,-128z_{14}^{2}\,-128z_{14}z_{23}\,-128z_{14}z_{28}\,-400z_{23}\,-32z_{23}^{2}$\\
 & $-64z_{23}z_{28}\,-432z_{28}\,-32z_{28}^{2}\,-2z_{29}\,-2z_{29}^{2}\,+z_{30}$\\[5pt]
$Q_{15}=$ & $-247\,-168z_{12}\,-32z_{12}^{2}\,-64z_{12}z_{13}\,-32z_{12}z_{27}\,-32z_{12}z_{28}$\\
 & $-188z_{13}\,-32z_{13}^{2}\,-32z_{13}z_{27}\,-32z_{13}z_{28}\,-2z_{23}\,-92z_{27}$\\
 & $-8z_{27}^{2}\,-16z_{27}z_{28}\,-84z_{28}\,-8z_{28}^{2}\,+z_{31}$\\[5pt]
$Q_{16}=$ & $-960\,-672z_{12}\,-128z_{12}^{2}\,-256z_{12}z_{13}\,-128z_{12}z_{27}\,-128z_{12}z_{28}$\\
 & $-736z_{13}\,-128z_{13}^{2}\,-128z_{13}z_{27}\,-128z_{13}z_{28}\,-368z_{27}\,-32z_{27}^{2}$\\
 & $-64z_{27}z_{28}\,-336z_{28}\,-32z_{28}^{2}\,-2z_{31}\,-2z_{31}^{2}\,+z_{32}$\\[5pt]
$Q_{17}=$ & $-276\,-30z_{12}\,-2z_{12}^{2}\,-16z_{12}z_{14}\,-8z_{12}z_{27}\,-30z_{13}$\\
 & $-2z_{13}^{2}\,-16z_{13}z_{14}\,-8z_{13}z_{28}\,-272z_{14}\,-64z_{14}^{2}\,-32z_{14}z_{27}$\\
 & $-32z_{14}z_{28}\,-68z_{27}\,-8z_{27}^{2}\,-68z_{28}\,-8z_{28}^{2}\,+z_{33}$\\[5pt]
$Q_{18}=$ & $-552\,-120z_{13}\,-8z_{13}^{2}\,-64z_{13}z_{14}\,-32z_{13}z_{28}\,-544z_{14}$\\
 & $-128z_{14}^{2}\,-128z_{14}z_{28}\,-272z_{28}\,-32z_{28}^{2}\,-2z_{33}\,-2z_{33}^{2}$\\
 & $+z_{34}$\\[5pt]
$Q_{19}=$ & $-247\,-4z_{12}\,-352z_{13}\,-128z_{13}^{2}\,-64z_{13}z_{23}\,-64z_{13}z_{27}$\\
 & $-92z_{23}\,-8z_{23}^{2}\,-16z_{23}z_{27}\,-84z_{27}\,-8z_{27}^{2}\,-2z_{28}$\\
 & $+z_{35}$\\[5pt]
$Q_{20}=$ & $-960\,-1408z_{13}\,-512z_{13}^{2}\,-256z_{13}z_{23}\,-256z_{13}z_{27}\,-368z_{23}$\\
 & $-32z_{23}^{2}\,-64z_{23}z_{27}\,-336z_{27}\,-32z_{27}^{2}\,-2z_{35}\,-2z_{35}^{2}$\\
 & $+z_{36}$\\[5pt]
$Q_{21}=$ & $-56\,-60z_{13}\,-16z_{13}^{2}\,-16z_{13}z_{23}\,-30z_{23}\,-4z_{23}^{2}$\\
 & $+z_{37}$\\[5pt]
$Q_{22}=$ & $-30\,-11z_{12}\,-z_{12}^{2}\,-4z_{12}z_{37}\,-22z_{37}\,-4z_{37}^{2}$\\
 & $+z_{38}$\\[5pt]
$Q_{23}=$ & $-380\,-312z_{37}\,-64z_{37}^{2}\,-32z_{37}z_{38}\,-78z_{38}\,-4z_{38}^{2}$\\
 & $+z_{39}$\\[5pt]
$Q_{24}=$ & $-30\,-11z_{13}\,-z_{13}^{2}\,-4z_{13}z_{39}\,-22z_{39}\,-4z_{39}^{2}$\\
 & $+z_{40}$\\[5pt]
$Q_{25}=$ & $-349\,-420z_{13}\,-128z_{13}^{2}\,-128z_{13}z_{23}\,-208z_{23}\,-32z_{23}^{2}$\\
 & $-2z_{27}\,+z_{41}$\\[5pt]
$Q_{26}=$ & $-1368\,-1664z_{13}\,-512z_{13}^{2}\,-512z_{13}z_{23}\,-832z_{23}\,-128z_{23}^{2}$\\
 & $-2z_{41}\,-2z_{41}^{2}\,+z_{42}$\\[5pt]
$Q_{27}=$ & $-56\,-60z_{12}\,-16z_{12}^{2}\,-16z_{12}z_{28}\,-8z_{12}z_{42}\,-30z_{28}$\\
 & $-4z_{28}^{2}\,-4z_{28}z_{42}\,-15z_{42}\,-z_{42}^{2}\,+z_{43}$\\[5pt]
$Q_{28}=$ & $-6\,-5z_{14}\,-z_{14}^{2}\,-2z_{14}z_{15}\,-5z_{15}\,-z_{15}^{2}$\\
 & $+z_{44}$\\[5pt]
$Q_{29}=$ & $-224\,-26z_{12}\,-2z_{12}^{2}\,-16z_{12}z_{15}\,-8z_{12}z_{44}\,-288z_{14}$\\
 & $-128z_{14}^{2}\,-64z_{14}z_{27}\,-64z_{14}z_{28}\,-88z_{15}\,-32z_{15}^{2}\,-32z_{15}z_{44}$\\
 & $-76z_{27}\,-8z_{27}^{2}\,-16z_{27}z_{28}\,-68z_{28}\,-8z_{28}^{2}\,-44z_{44}$\\
 & $-8z_{44}^{2}\,+z_{45}$\\[5pt]
$Q_{30}=$ & $-640\,-1152z_{14}\,-512z_{14}^{2}\,-256z_{14}z_{27}\,-256z_{14}z_{28}\,-304z_{27}$\\
 & $-32z_{27}^{2}\,-64z_{27}z_{28}\,-272z_{28}\,-32z_{28}^{2}\,-2z_{45}\,-2z_{45}^{2}$\\
 & $+z_{46}$\\[5pt]
$Q_{31}=$ & $-241\,-168z_{12}\,-32z_{12}^{2}\,-64z_{12}z_{13}\,-32z_{12}z_{27}\,-32z_{12}z_{28}$\\
 & $-184z_{13}\,-32z_{13}^{2}\,-32z_{13}z_{27}\,-32z_{13}z_{28}\,-z_{14}\,-92z_{27}$\\
 & $-8z_{27}^{2}\,-16z_{27}z_{28}\,-84z_{28}\,-8z_{28}^{2}\,+z_{47}$\\[5pt]
$Q_{32}=$ & $-42\,-52z_{14}\,-16z_{14}^{2}\,-16z_{14}z_{27}\,-26z_{27}\,-4z_{27}^{2}$\\
 & $+z_{48}$\\[5pt]
$Q_{33}=$ & $-30\,-11z_{14}\,-z_{14}^{2}\,-4z_{14}z_{48}\,-22z_{48}\,-4z_{48}^{2}$\\
 & $+z_{49}$\\[5pt]
$Q_{34}=$ & $-306\,-35z_{13}\,-z_{13}^{2}\,-16z_{13}z_{48}\,-4z_{13}z_{49}\,-280z_{48}$\\
 & $-64z_{48}^{2}\,-32z_{48}z_{49}\,-70z_{49}\,-4z_{49}^{2}\,+z_{50}$\\[5pt]
$Q_{35}=$ & $-4556\,-4320z_{48}\,-1024z_{48}^{2}\,-512z_{48}z_{49}\,-128z_{48}z_{50}\,-1080z_{49}$\\
 & $-64z_{49}^{2}\,-32z_{49}z_{50}\,-270z_{50}\,-4z_{50}^{2}\,+z_{51}$\\[5pt]
$Q_{36}=$ & $-30\,-11z_{14}\,-z_{14}^{2}\,-4z_{14}z_{51}\,-22z_{51}\,-4z_{51}^{2}$\\
 & $+z_{52}$\\[5pt]
$Q_{37}=$ & $-306\,-35z_{13}\,-z_{13}^{2}\,-16z_{13}z_{51}\,-4z_{13}z_{52}\,-280z_{51}$\\
 & $-64z_{51}^{2}\,-32z_{51}z_{52}\,-70z_{52}\,-4z_{52}^{2}\,+z_{53}$\\[5pt]
$Q_{38}=$ & $-1028\,-672z_{12}\,-128z_{12}^{2}\,-256z_{12}z_{13}\,-128z_{12}z_{27}\,-128z_{12}z_{28}$\\
 & $-736z_{13}\,-128z_{13}^{2}\,-128z_{13}z_{27}\,-128z_{13}z_{28}\,-368z_{27}\,-32z_{27}^{2}$\\
 & $-64z_{27}z_{28}\,-336z_{28}\,-32z_{28}^{2}\,-2z_{47}\,-2z_{47}^{2}\,-32z_{51}$\\
 & $-8z_{52}\,-2z_{53}\,+z_{54}$\\[5pt]
$Q_{39}=$ & $-264\,-128z_{51}\,-32z_{52}\,-8z_{53}\,-2z_{54}\,-2z_{54}^{2}$\\
 & $+z_{55}$\\[5pt]
$Q_{40}=$ & $-2\,-3z_{12}\,-z_{12}^{2}\,-2z_{12}z_{55}\,-3z_{55}\,-z_{55}^{2}$\\
 & $+z_{56}$\\[5pt]
$Q_{41}=$ & $-12\,-28z_{55}\,-16z_{55}^{2}\,-16z_{55}z_{56}\,-14z_{56}\,-4z_{56}^{2}$\\
 & $+z_{57}$\\[5pt]
$Q_{42}=$ & $-30\,-11z_{14}\,-z_{14}^{2}\,-4z_{14}z_{57}\,-22z_{57}\,-4z_{57}^{2}$\\
 & $+z_{58}$\\[5pt]
$Q_{43}=$ & $-306\,-35z_{12}\,-z_{12}^{2}\,-16z_{12}z_{57}\,-4z_{12}z_{58}\,-280z_{57}$\\
 & $-64z_{57}^{2}\,-32z_{57}z_{58}\,-70z_{58}\,-4z_{58}^{2}\,+z_{59}$\\[5pt]
$Q_{44}=$ & $-4556\,-4320z_{57}\,-1024z_{57}^{2}\,-512z_{57}z_{58}\,-128z_{57}z_{59}\,-1080z_{58}$\\
 & $-64z_{58}^{2}\,-32z_{58}z_{59}\,-270z_{59}\,-4z_{59}^{2}\,+z_{60}$\\[5pt]
$Q_{45}=$ & $-90\,-76z_{14}\,-16z_{14}^{2}\,-16z_{14}z_{27}\,-16z_{14}z_{60}\,-38z_{27}$\\
 & $-4z_{27}^{2}\,-8z_{27}z_{60}\,-38z_{60}\,-4z_{60}^{2}\,+z_{61}$\\[5pt]
$Q_{46}=$ & $-4556\,-4320z_{51}\,-1024z_{51}^{2}\,-512z_{51}z_{52}\,-128z_{51}z_{53}\,-1080z_{52}$\\
 & $-64z_{52}^{2}\,-32z_{52}z_{53}\,-270z_{53}\,-4z_{53}^{2}\,+z_{62}$\\[5pt]
$Q_{47}=$ & $-682\,-70z_{13}\,-2z_{13}^{2}\,-32z_{13}z_{60}\,-8z_{13}z_{61}\,-592z_{60}$\\
 & $-128z_{60}^{2}\,-64z_{60}z_{61}\,-148z_{61}\,-8z_{61}^{2}\,-2z_{62}\,+z_{63}$\\[5pt]
$Q_{48}=$ & $-z_{63}\,-z_{63}^{2}\,+z_{64}$\\[5pt]
$Q_{49}=$ & $-42\,-52z_{13}\,-16z_{13}^{2}\,-8z_{13}z_{14}\,-16z_{13}z_{23}\,-13z_{14}$\\
 & $-z_{14}^{2}\,-4z_{14}z_{23}\,-26z_{23}\,-4z_{23}^{2}\,+z_{65}$\\[5pt]
$Q_{50}=$ & $-2278\,-2264z_{13}\,-800z_{13}^{2}\,-800z_{13}z_{23}\,-160z_{13}z_{65}\,-70z_{14}$\\
 & $-2z_{14}^{2}\,-32z_{14}z_{62}\,-8z_{14}z_{64}\,-1132z_{23}\,-200z_{23}^{2}\,-80z_{23}z_{65}$\\
 & $-592z_{62}\,-128z_{62}^{2}\,-64z_{62}z_{64}\,-148z_{64}\,-8z_{64}^{2}\,-228z_{65}$\\
 & $-8z_{65}^{2}\,+z_{66}$\\[5pt]
$Q_{51}=$ & $-2712\,-280z_{14}\,-8z_{14}^{2}\,-128z_{14}z_{62}\,-32z_{14}z_{64}\,-2368z_{62}$\\
 & $-512z_{62}^{2}\,-256z_{62}z_{64}\,-592z_{64}\,-32z_{64}^{2}\,-2z_{66}\,-2z_{66}^{2}$\\
 & $+z_{67}$\\[5pt]
$Q_{52}=$ & $-2\,-3z_{67}\,-z_{67}^{2}\,+z_{68}$\\[5pt]
$Q_{53}=$ & $-z_{2}\,-z_{2}z_{17}\,+z_{69}$\\[5pt]
$Q_{54}=$ & $-z_{16}^{2}\,+2z_{16}z_{17}\,-z_{17}^{2}\,+z_{70}$\\[5pt]
$Q_{55}=$ & $-2z_{26}z_{69}\,-4z_{69}\,+z_{71}$\\[5pt]
$Q_{56}=$ & $-z_{16}^{2}\,+2z_{16}z_{17}\,+2z_{16}z_{71}\,-z_{17}^{2}\,-2z_{17}z_{71}\,-z_{71}^{2}$\\
 & $+z_{72}$\\[5pt]
$Q_{57}=$ & $-z_{30}z_{69}\,+z_{73}$\\[5pt]
$Q_{58}=$ & $-z_{16}^{2}\,+2z_{16}z_{17}\,+2z_{16}z_{73}\,-z_{17}^{2}\,-2z_{17}z_{73}\,-z_{73}^{2}$\\
 & $+z_{74}$\\[5pt]
$Q_{59}=$ & $-z_{32}z_{69}\,+z_{75}$\\[5pt]
$Q_{60}=$ & $-z_{16}^{2}\,+2z_{16}z_{17}\,+2z_{16}z_{75}\,-z_{17}^{2}\,-2z_{17}z_{75}\,-z_{75}^{2}$\\
 & $+z_{76}$\\[5pt]
$Q_{61}=$ & $-z_{34}z_{69}\,+z_{77}$\\[5pt]
$Q_{62}=$ & $-z_{16}^{2}\,+2z_{16}z_{17}\,+2z_{16}z_{77}\,-z_{17}^{2}\,-2z_{17}z_{77}\,-z_{77}^{2}$\\
 & $+z_{78}$\\[5pt]
$Q_{63}=$ & $-z_{36}z_{69}\,+z_{79}$\\[5pt]
$Q_{64}=$ & $-z_{16}^{2}\,+2z_{16}z_{17}\,+2z_{16}z_{79}\,-z_{17}^{2}\,-2z_{17}z_{79}\,-z_{79}^{2}$\\
 & $+z_{80}$\\[5pt]
$Q_{65}=$ & $-4z_{13}z_{69}\,-2z_{40}z_{69}\,-4z_{69}\,+z_{81}$\\[5pt]
$Q_{66}=$ & $-z_{16}^{2}\,+2z_{16}z_{17}\,+2z_{16}z_{81}\,-z_{17}^{2}\,-2z_{17}z_{81}\,-z_{81}^{2}$\\
 & $+z_{82}$\\[5pt]
$Q_{67}=$ & $-4z_{42}z_{69}\,-2z_{43}z_{69}\,+z_{83}$\\[5pt]
$Q_{68}=$ & $-z_{16}^{2}\,+2z_{16}z_{17}\,+2z_{16}z_{83}\,-z_{17}^{2}\,-2z_{17}z_{83}\,-z_{83}^{2}$\\
 & $+z_{84}$\\[5pt]
$Q_{69}=$ & $-z_{46}z_{69}\,+z_{85}$\\[5pt]
$Q_{70}=$ & $-z_{16}^{2}\,+2z_{16}z_{17}\,+2z_{16}z_{85}\,-z_{17}^{2}\,-2z_{17}z_{85}\,-z_{85}^{2}$\\
 & $+z_{86}$\\[5pt]
$Q_{71}=$ & $-z_{12}^{2}\,+2z_{12}z_{14}\,-z_{14}^{2}\,+z_{87}$\\[5pt]
$Q_{72}=$ & $-z_{13}^{2}\,+2z_{13}z_{14}\,-z_{14}^{2}\,+z_{88}$\\[5pt]
$Q_{73}=$ & $-z_{87}z_{88}\,+z_{89}$\\[5pt]
$Q_{74}=$ & $-4\,-4z_{15}\,-z_{15}^{2}\,+z_{90}$\\[5pt]
$Q_{75}=$ & $-z_{89}^{2}\,+2z_{89}z_{90}\,-z_{90}^{2}\,+z_{91}$\\[5pt]
$Q_{76}=$ & $-2z_{68}z_{69}\,-4z_{69}\,+z_{92}$\\[5pt]
$Q_{77}=$ & $-z_{16}^{2}\,+2z_{16}z_{17}\,+2z_{16}z_{92}\,-z_{17}^{2}\,-2z_{17}z_{92}\,-z_{92}^{2}$\\
 & $+z_{93}$\\[5pt]
$Q_{78}=$ & $-z_{2}\,-z_{2}z_{13}\,+z_{94}$\\[5pt]
$Q_{79}=$ & $-1\,-z_{17}\,-z_{17}z_{94}\,-z_{94}\,+z_{95}$\\[5pt]
$Q_{80}=$ & $-z_{2}\,-z_{2}z_{12}\,+z_{96}$\\[5pt]
$Q_{81}=$ & $-4z_{2}\,-4z_{2}z_{13}\,-2z_{2}z_{23}\,+z_{97}$\\[5pt]
$Q_{82}=$ & $-1\,-z_{16}\,-z_{16}z_{96}\,-z_{96}\,+z_{98}$\\[5pt]
$Q_{83}=$ & $-z_{97}z_{98}\,-z_{98}\,+z_{99}$\\[5pt]
$Q_{84}=$ & $-z_{95}^{2}\,+2z_{95}z_{99}\,-z_{99}^{2}\,+z_{100}$\\[5pt]
$Q_{85}=$ & $-6z_{2}\,-4z_{2}z_{13}\,-2z_{2}z_{23}\,+z_{101}$\\[5pt]
$Q_{86}=$ & $-1\,-z_{17}\,-z_{17}z_{101}\,-z_{101}\,+z_{102}$\\[5pt]
$Q_{87}=$ & $-1\,-z_{16}\,-z_{16}z_{94}\,-z_{94}\,+z_{103}$\\[5pt]
$Q_{88}=$ & $-z_{102}^{2}\,+2z_{102}z_{103}\,-z_{103}^{2}\,+z_{104}$\\[5pt]
$Q_{89}=$ & $-z_{70}z_{72}\,+z_{105}$\\[5pt]
$Q_{90}=$ & $-z_{74}z_{105}\,+z_{106}$\\[5pt]
$Q_{91}=$ & $-z_{76}z_{106}\,+z_{107}$\\[5pt]
$Q_{92}=$ & $-z_{78}z_{107}\,+z_{108}$\\[5pt]
$Q_{93}=$ & $-z_{80}z_{108}\,+z_{109}$\\[5pt]
$Q_{94}=$ & $-z_{82}z_{109}\,+z_{110}$\\[5pt]
$Q_{95}=$ & $-z_{84}z_{110}\,+z_{111}$\\[5pt]
$Q_{96}=$ & $-z_{86}z_{111}\,+z_{112}$\\[5pt]
$Q_{97}=$ & $-z_{91}z_{112}\,-z_{93}z_{112}\,+z_{113}$\\[5pt]
$Q_{98}=$ & $-z_{100}z_{113}\,+z_{114}$\\[5pt]
$Q_{99}=$ & $-2\,-2z_{2}\,-z_{2}z_{20}\,-2z_{9}\,-z_{9}z_{20}\,-z_{20}$\\
 & $+z_{115}$\\[5pt]
$Q_{100}=$ & $-1\,-z_{18}\,-z_{18}z_{115}\,-z_{115}\,+z_{116}$\\[5pt]
$Q_{101}=$ & $-2\,-z_{2}\,-z_{2}z_{19}\,-z_{9}\,-z_{9}z_{19}\,-2z_{19}$\\
 & $-z_{19}z_{115}\,-z_{115}\,+z_{117}$\\[5pt]
$Q_{102}=$ & $2\,-4z_{1}\,+2z_{1}^{2}\,-2z_{1}z_{16}\,-2z_{1}z_{17}\,-2z_{1}z_{116}$\\
 & $-2z_{1}z_{117}\,+2z_{16}\,+z_{16}^{2}\,+2z_{16}z_{116}\,+2z_{17}\,+z_{17}^{2}$\\
 & $+2z_{17}z_{117}\,-z_{21}z_{118}\,+z_{104}z_{114}\,+2z_{116}\,+z_{116}^{2}\,+2z_{117}$\\
 & $+z_{117}^{2}$\\[5pt]
$Q_{103}=$ & $3\,+z_{4}\,-z_{21}z_{119}\,+z_{173}$\\[5pt]
$Q_{104}=$ & $-2\,-z_{4}\,-z_{4}z_{171}\,-z_{4}z_{172}\,-3z_{171}\,-z_{171}z_{173}$\\
 & $-3z_{172}\,-z_{172}z_{173}\,-z_{173}\,+z_{178}\,+z_{179}$\\[5pt]
$Q_{105}=$ & $-3\,-z_{4}\,-z_{4}z_{178}\,-z_{4}z_{179}\,+z_{120}\,-z_{173}$\\
 & $-z_{173}z_{178}\,-z_{173}z_{179}\,-3z_{178}\,-3z_{179}$\\[5pt]
$Q_{106}=$ & $-z_{120}\,-z_{120}z_{203}\,-z_{120}z_{204}\,+z_{121}$\\[5pt]
$Q_{107}=$ & $-z_{121}\,-z_{121}z_{216}\,-z_{121}z_{217}\,+z_{122}$\\[5pt]
$Q_{108}=$ & $-1\,-z_{125}\,-z_{125}z_{216}\,-z_{125}z_{217}\,-z_{125}z_{240}\,+z_{126}$\\
 & $-z_{216}\,-z_{217}\,-z_{240}$\\[5pt]
$Q_{109}=$ & $-2z_{126}\,-z_{126}z_{216}\,-z_{126}z_{217}\,-z_{126}z_{240}\,+z_{127}$\\[5pt]
$Q_{110}=$ & $-1\,+z_{123}^{2}\,+z_{126}^{2}\,-z_{127}^{2}$\\[5pt]
$Q_{111}=$ & $-17\,-8z_{4}\,-z_{4}^{2}\,+2z_{4}z_{123}\,+2z_{4}z_{124}\,-2z_{4}z_{173}$\\
 & $-z_{122}z_{128}\,+8z_{123}\,+2z_{123}z_{173}\,+8z_{124}\,+2z_{124}z_{173}\,-8z_{173}$\\
 & $-z_{173}^{2}$\\[5pt]
$Q_{112}=$ & $-8\,-8z_{134}\,-4z_{134}z_{145}\,-4z_{134}z_{150}\,-4z_{134}z_{160}\,-4z_{134}z_{164}$\\
 & $-4z_{134}z_{168}\,-4z_{134}z_{175}\,-4z_{134}z_{182}\,-4z_{134}z_{188}\,-4z_{134}z_{194}\,-4z_{134}z_{200}$\\
 & $-4z_{134}z_{207}\,-4z_{134}z_{213}\,-4z_{134}z_{216}\,-4z_{134}z_{217}\,-4z_{134}z_{220}\,-4z_{134}z_{226}$\\
 & $-4z_{134}z_{233}\,-4z_{134}z_{238}\,-4z_{134}z_{240}\,+z_{135}\,-4z_{145}\,-4z_{150}$\\
 & $-4z_{160}\,-4z_{164}\,-4z_{168}\,-4z_{175}\,-4z_{182}\,-4z_{188}$\\
 & $-4z_{194}\,-4z_{200}\,-4z_{207}\,-4z_{213}\,-4z_{216}\,-4z_{217}$\\
 & $-4z_{220}\,-4z_{226}\,-4z_{233}\,-4z_{238}\,-4z_{240}$\\[5pt]
$Q_{113}=$ & $-2z_{123}\,-z_{123}z_{145}\,-z_{123}z_{150}\,-z_{123}z_{160}\,-z_{123}z_{164}\,-z_{123}z_{168}$\\
 & $-z_{123}z_{175}\,-z_{123}z_{182}\,-z_{123}z_{188}\,-z_{123}z_{194}\,-z_{123}z_{200}\,-z_{123}z_{207}$\\
 & $-z_{123}z_{213}\,-z_{123}z_{216}\,-z_{123}z_{217}\,-z_{123}z_{220}\,-z_{123}z_{226}\,-z_{123}z_{233}$\\
 & $-z_{123}z_{238}\,-z_{123}z_{240}\,-2z_{124}\,-z_{124}z_{145}\,-z_{124}z_{150}\,-z_{124}z_{160}$\\
 & $-z_{124}z_{164}\,-z_{124}z_{168}\,-z_{124}z_{175}\,-z_{124}z_{182}\,-z_{124}z_{188}\,-z_{124}z_{194}$\\
 & $-z_{124}z_{200}\,-z_{124}z_{207}\,-z_{124}z_{213}\,-z_{124}z_{216}\,-z_{124}z_{217}\,-z_{124}z_{220}$\\
 & $-z_{124}z_{226}\,-z_{124}z_{233}\,-z_{124}z_{238}\,-z_{124}z_{240}\,+z_{136}$\\[5pt]
$Q_{114}=$ & $3\,+z_{129}^{2}\,-z_{136}^{2}\,+4z_{145}\,+z_{145}^{2}\,+2z_{145}z_{150}$\\
 & $+2z_{145}z_{160}\,+2z_{145}z_{164}\,+2z_{145}z_{168}\,+2z_{145}z_{175}\,+2z_{145}z_{182}\,+2z_{145}z_{188}$\\
 & $+2z_{145}z_{194}\,+2z_{145}z_{200}\,+2z_{145}z_{207}\,+2z_{145}z_{213}\,+2z_{145}z_{216}\,+2z_{145}z_{217}$\\
 & $+2z_{145}z_{220}\,+2z_{145}z_{226}\,+2z_{145}z_{233}\,+2z_{145}z_{238}\,+2z_{145}z_{240}\,+4z_{150}$\\
 & $+z_{150}^{2}\,+2z_{150}z_{160}\,+2z_{150}z_{164}\,+2z_{150}z_{168}\,+2z_{150}z_{175}\,+2z_{150}z_{182}$\\
 & $+2z_{150}z_{188}\,+2z_{150}z_{194}\,+2z_{150}z_{200}\,+2z_{150}z_{207}\,+2z_{150}z_{213}\,+2z_{150}z_{216}$\\
 & $+2z_{150}z_{217}\,+2z_{150}z_{220}\,+2z_{150}z_{226}\,+2z_{150}z_{233}\,+2z_{150}z_{238}\,+2z_{150}z_{240}$\\
 & $+4z_{160}\,+z_{160}^{2}\,+2z_{160}z_{164}\,+2z_{160}z_{168}\,+2z_{160}z_{175}\,+2z_{160}z_{182}$\\
 & $+2z_{160}z_{188}\,+2z_{160}z_{194}\,+2z_{160}z_{200}\,+2z_{160}z_{207}\,+2z_{160}z_{213}\,+2z_{160}z_{216}$\\
 & $+2z_{160}z_{217}\,+2z_{160}z_{220}\,+2z_{160}z_{226}\,+2z_{160}z_{233}\,+2z_{160}z_{238}\,+2z_{160}z_{240}$\\
 & $+4z_{164}\,+z_{164}^{2}\,+2z_{164}z_{168}\,+2z_{164}z_{175}\,+2z_{164}z_{182}\,+2z_{164}z_{188}$\\
 & $+2z_{164}z_{194}\,+2z_{164}z_{200}\,+2z_{164}z_{207}\,+2z_{164}z_{213}\,+2z_{164}z_{216}\,+2z_{164}z_{217}$\\
 & $+2z_{164}z_{220}\,+2z_{164}z_{226}\,+2z_{164}z_{233}\,+2z_{164}z_{238}\,+2z_{164}z_{240}\,+4z_{168}$\\
 & $+z_{168}^{2}\,+2z_{168}z_{175}\,+2z_{168}z_{182}\,+2z_{168}z_{188}\,+2z_{168}z_{194}\,+2z_{168}z_{200}$\\
 & $+2z_{168}z_{207}\,+2z_{168}z_{213}\,+2z_{168}z_{216}\,+2z_{168}z_{217}\,+2z_{168}z_{220}\,+2z_{168}z_{226}$\\
 & $+2z_{168}z_{233}\,+2z_{168}z_{238}\,+2z_{168}z_{240}\,+4z_{175}\,+z_{175}^{2}\,+2z_{175}z_{182}$\\
 & $+2z_{175}z_{188}\,+2z_{175}z_{194}\,+2z_{175}z_{200}\,+2z_{175}z_{207}\,+2z_{175}z_{213}\,+2z_{175}z_{216}$\\
 & $+2z_{175}z_{217}\,+2z_{175}z_{220}\,+2z_{175}z_{226}\,+2z_{175}z_{233}\,+2z_{175}z_{238}\,+2z_{175}z_{240}$\\
 & $+4z_{182}\,+z_{182}^{2}\,+2z_{182}z_{188}\,+2z_{182}z_{194}\,+2z_{182}z_{200}\,+2z_{182}z_{207}$\\
 & $+2z_{182}z_{213}\,+2z_{182}z_{216}\,+2z_{182}z_{217}\,+2z_{182}z_{220}\,+2z_{182}z_{226}\,+2z_{182}z_{233}$\\
 & $+2z_{182}z_{238}\,+2z_{182}z_{240}\,+4z_{188}\,+z_{188}^{2}\,+2z_{188}z_{194}\,+2z_{188}z_{200}$\\
 & $+2z_{188}z_{207}\,+2z_{188}z_{213}\,+2z_{188}z_{216}\,+2z_{188}z_{217}\,+2z_{188}z_{220}\,+2z_{188}z_{226}$\\
 & $+2z_{188}z_{233}\,+2z_{188}z_{238}\,+2z_{188}z_{240}\,+4z_{194}\,+z_{194}^{2}\,+2z_{194}z_{200}$\\
 & $+2z_{194}z_{207}\,+2z_{194}z_{213}\,+2z_{194}z_{216}\,+2z_{194}z_{217}\,+2z_{194}z_{220}\,+2z_{194}z_{226}$\\
 & $+2z_{194}z_{233}\,+2z_{194}z_{238}\,+2z_{194}z_{240}\,+4z_{200}\,+z_{200}^{2}\,+2z_{200}z_{207}$\\
 & $+2z_{200}z_{213}\,+2z_{200}z_{216}\,+2z_{200}z_{217}\,+2z_{200}z_{220}\,+2z_{200}z_{226}\,+2z_{200}z_{233}$\\
 & $+2z_{200}z_{238}\,+2z_{200}z_{240}\,+4z_{207}\,+z_{207}^{2}\,+2z_{207}z_{213}\,+2z_{207}z_{216}$\\
 & $+2z_{207}z_{217}\,+2z_{207}z_{220}\,+2z_{207}z_{226}\,+2z_{207}z_{233}\,+2z_{207}z_{238}\,+2z_{207}z_{240}$\\
 & $+4z_{213}\,+z_{213}^{2}\,+2z_{213}z_{216}\,+2z_{213}z_{217}\,+2z_{213}z_{220}\,+2z_{213}z_{226}$\\
 & $+2z_{213}z_{233}\,+2z_{213}z_{238}\,+2z_{213}z_{240}\,+4z_{216}\,+z_{216}^{2}\,+2z_{216}z_{217}$\\
 & $+2z_{216}z_{220}\,+2z_{216}z_{226}\,+2z_{216}z_{233}\,+2z_{216}z_{238}\,+2z_{216}z_{240}\,+4z_{217}$\\
 & $+z_{217}^{2}\,+2z_{217}z_{220}\,+2z_{217}z_{226}\,+2z_{217}z_{233}\,+2z_{217}z_{238}\,+2z_{217}z_{240}$\\
 & $+4z_{220}\,+z_{220}^{2}\,+2z_{220}z_{226}\,+2z_{220}z_{233}\,+2z_{220}z_{238}\,+2z_{220}z_{240}$\\
 & $+4z_{226}\,+z_{226}^{2}\,+2z_{226}z_{233}\,+2z_{226}z_{238}\,+2z_{226}z_{240}\,+4z_{233}$\\
 & $+z_{233}^{2}\,+2z_{233}z_{238}\,+2z_{233}z_{240}\,+4z_{238}\,+z_{238}^{2}\,+2z_{238}z_{240}$\\
 & $+4z_{240}\,+z_{240}^{2}$\\[5pt]
$Q_{115}=$ & $-z_{123}z_{131}\,-z_{124}z_{131}\,+z_{137}$\\[5pt]
$Q_{116}=$ & $-1\,+z_{130}^{2}\,+z_{131}^{2}\,-z_{137}^{2}$\\[5pt]
$Q_{117}=$ & $-z_{133}\,-z_{133}z_{135}\,+z_{138}$\\[5pt]
$Q_{118}=$ & $-1\,+z_{132}^{2}\,+z_{133}^{2}\,-z_{138}^{2}$\\[5pt]
$Q_{119}=$ & $-4\,+z_{140}\,-4z_{145}\,-z_{145}^{2}\,-2z_{145}z_{150}\,-2z_{145}z_{160}$\\
 & $-2z_{145}z_{164}\,-2z_{145}z_{168}\,-2z_{145}z_{175}\,-2z_{145}z_{182}\,-2z_{145}z_{188}\,-2z_{145}z_{194}$\\
 & $-2z_{145}z_{200}\,-2z_{145}z_{207}\,-2z_{145}z_{213}\,-2z_{145}z_{216}\,-2z_{145}z_{217}\,-2z_{145}z_{220}$\\
 & $-2z_{145}z_{226}\,-2z_{145}z_{233}\,-2z_{145}z_{238}\,-2z_{145}z_{240}\,-4z_{150}\,-z_{150}^{2}$\\
 & $-2z_{150}z_{160}\,-2z_{150}z_{164}\,-2z_{150}z_{168}\,-2z_{150}z_{175}\,-2z_{150}z_{182}\,-2z_{150}z_{188}$\\
 & $-2z_{150}z_{194}\,-2z_{150}z_{200}\,-2z_{150}z_{207}\,-2z_{150}z_{213}\,-2z_{150}z_{216}\,-2z_{150}z_{217}$\\
 & $-2z_{150}z_{220}\,-2z_{150}z_{226}\,-2z_{150}z_{233}\,-2z_{150}z_{238}\,-2z_{150}z_{240}\,-4z_{160}$\\
 & $-z_{160}^{2}\,-2z_{160}z_{164}\,-2z_{160}z_{168}\,-2z_{160}z_{175}\,-2z_{160}z_{182}\,-2z_{160}z_{188}$\\
 & $-2z_{160}z_{194}\,-2z_{160}z_{200}\,-2z_{160}z_{207}\,-2z_{160}z_{213}\,-2z_{160}z_{216}\,-2z_{160}z_{217}$\\
 & $-2z_{160}z_{220}\,-2z_{160}z_{226}\,-2z_{160}z_{233}\,-2z_{160}z_{238}\,-2z_{160}z_{240}\,-4z_{164}$\\
 & $-z_{164}^{2}\,-2z_{164}z_{168}\,-2z_{164}z_{175}\,-2z_{164}z_{182}\,-2z_{164}z_{188}\,-2z_{164}z_{194}$\\
 & $-2z_{164}z_{200}\,-2z_{164}z_{207}\,-2z_{164}z_{213}\,-2z_{164}z_{216}\,-2z_{164}z_{217}\,-2z_{164}z_{220}$\\
 & $-2z_{164}z_{226}\,-2z_{164}z_{233}\,-2z_{164}z_{238}\,-2z_{164}z_{240}\,-4z_{168}\,-z_{168}^{2}$\\
 & $-2z_{168}z_{175}\,-2z_{168}z_{182}\,-2z_{168}z_{188}\,-2z_{168}z_{194}\,-2z_{168}z_{200}\,-2z_{168}z_{207}$\\
 & $-2z_{168}z_{213}\,-2z_{168}z_{216}\,-2z_{168}z_{217}\,-2z_{168}z_{220}\,-2z_{168}z_{226}\,-2z_{168}z_{233}$\\
 & $-2z_{168}z_{238}\,-2z_{168}z_{240}\,-4z_{175}\,-z_{175}^{2}\,-2z_{175}z_{182}\,-2z_{175}z_{188}$\\
 & $-2z_{175}z_{194}\,-2z_{175}z_{200}\,-2z_{175}z_{207}\,-2z_{175}z_{213}\,-2z_{175}z_{216}\,-2z_{175}z_{217}$\\
 & $-2z_{175}z_{220}\,-2z_{175}z_{226}\,-2z_{175}z_{233}\,-2z_{175}z_{238}\,-2z_{175}z_{240}\,-4z_{182}$\\
 & $-z_{182}^{2}\,-2z_{182}z_{188}\,-2z_{182}z_{194}\,-2z_{182}z_{200}\,-2z_{182}z_{207}\,-2z_{182}z_{213}$\\
 & $-2z_{182}z_{216}\,-2z_{182}z_{217}\,-2z_{182}z_{220}\,-2z_{182}z_{226}\,-2z_{182}z_{233}\,-2z_{182}z_{238}$\\
 & $-2z_{182}z_{240}\,-4z_{188}\,-z_{188}^{2}\,-2z_{188}z_{194}\,-2z_{188}z_{200}\,-2z_{188}z_{207}$\\
 & $-2z_{188}z_{213}\,-2z_{188}z_{216}\,-2z_{188}z_{217}\,-2z_{188}z_{220}\,-2z_{188}z_{226}\,-2z_{188}z_{233}$\\
 & $-2z_{188}z_{238}\,-2z_{188}z_{240}\,-4z_{194}\,-z_{194}^{2}\,-2z_{194}z_{200}\,-2z_{194}z_{207}$\\
 & $-2z_{194}z_{213}\,-2z_{194}z_{216}\,-2z_{194}z_{217}\,-2z_{194}z_{220}\,-2z_{194}z_{226}\,-2z_{194}z_{233}$\\
 & $-2z_{194}z_{238}\,-2z_{194}z_{240}\,-4z_{200}\,-z_{200}^{2}\,-2z_{200}z_{207}\,-2z_{200}z_{213}$\\
 & $-2z_{200}z_{216}\,-2z_{200}z_{217}\,-2z_{200}z_{220}\,-2z_{200}z_{226}\,-2z_{200}z_{233}\,-2z_{200}z_{238}$\\
 & $-2z_{200}z_{240}\,-4z_{207}\,-z_{207}^{2}\,-2z_{207}z_{213}\,-2z_{207}z_{216}\,-2z_{207}z_{217}$\\
 & $-2z_{207}z_{220}\,-2z_{207}z_{226}\,-2z_{207}z_{233}\,-2z_{207}z_{238}\,-2z_{207}z_{240}\,-4z_{213}$\\
 & $-z_{213}^{2}\,-2z_{213}z_{216}\,-2z_{213}z_{217}\,-2z_{213}z_{220}\,-2z_{213}z_{226}\,-2z_{213}z_{233}$\\
 & $-2z_{213}z_{238}\,-2z_{213}z_{240}\,-4z_{216}\,-z_{216}^{2}\,-2z_{216}z_{217}\,-2z_{216}z_{220}$\\
 & $-2z_{216}z_{226}\,-2z_{216}z_{233}\,-2z_{216}z_{238}\,-2z_{216}z_{240}\,-4z_{217}\,-z_{217}^{2}$\\
 & $-2z_{217}z_{220}\,-2z_{217}z_{226}\,-2z_{217}z_{233}\,-2z_{217}z_{238}\,-2z_{217}z_{240}\,-4z_{220}$\\
 & $-z_{220}^{2}\,-2z_{220}z_{226}\,-2z_{220}z_{233}\,-2z_{220}z_{238}\,-2z_{220}z_{240}\,-4z_{226}$\\
 & $-z_{226}^{2}\,-2z_{226}z_{233}\,-2z_{226}z_{238}\,-2z_{226}z_{240}\,-4z_{233}\,-z_{233}^{2}$\\
 & $-2z_{233}z_{238}\,-2z_{233}z_{240}\,-4z_{238}\,-z_{238}^{2}\,-2z_{238}z_{240}\,-4z_{240}$\\
 & $-z_{240}^{2}$\\[5pt]
$Q_{120}=$ & $z_{131}\,-z_{139}z_{140}\,-z_{140}$\\[5pt]
$Q_{121}=$ & $1\,-z_{123}\,-z_{124}\,-z_{130}z_{141}\,+z_{135}$\\[5pt]
$Q_{122}=$ & $-z_{129}\,-z_{130}z_{142}\,+z_{132}$\\[5pt]
$Q_{123}=$ & $-1\,+z_{133}\,-8z_{143}\,-4z_{143}z_{145}\,-4z_{143}z_{150}\,-4z_{143}z_{160}$\\
 & $-4z_{143}z_{164}\,-4z_{143}z_{168}\,-4z_{143}z_{175}\,-4z_{143}z_{182}\,-4z_{143}z_{188}\,-4z_{143}z_{194}$\\
 & $-4z_{143}z_{200}\,-4z_{143}z_{207}\,-4z_{143}z_{213}\,-4z_{143}z_{216}\,-4z_{143}z_{217}\,-4z_{143}z_{220}$\\
 & $-4z_{143}z_{226}\,-4z_{143}z_{233}\,-4z_{143}z_{238}\,-4z_{143}z_{240}\,-z_{216}\,-z_{217}$\\
 & $-z_{240}$\\[5pt]
$Q_{124}=$ & $-z_{123}z_{145}\,-z_{124}z_{145}\,+z_{146}$\\[5pt]
$Q_{125}=$ & $-1\,+z_{144}^{2}\,+z_{145}^{2}\,-z_{146}^{2}$\\[5pt]
$Q_{126}=$ & $-2\,-2z_{3}\,-z_{3}^{2}\,+2z_{3}z_{123}\,+2z_{3}z_{124}\,+2z_{123}$\\
 & $+2z_{124}\,-z_{147}$\\[5pt]
$Q_{127}=$ & $z_{3}z_{145}\,-z_{5}\,+z_{144}\,+z_{145}\,-z_{146}\,-z_{147}z_{148}$\\[5pt]
$Q_{128}=$ & $-z_{123}z_{150}\,-z_{124}z_{150}\,+z_{151}$\\[5pt]
$Q_{129}=$ & $-1\,+z_{149}^{2}\,+z_{150}^{2}\,-z_{151}^{2}$\\[5pt]
$Q_{130}=$ & $-5\,-4z_{4}\,-z_{4}^{2}\,+2z_{4}z_{123}\,+2z_{4}z_{124}\,+4z_{123}$\\
 & $+4z_{124}\,-z_{152}$\\[5pt]
$Q_{131}=$ & $z_{4}z_{150}\,-z_{21}\,+z_{149}\,+2z_{150}\,-z_{151}\,-z_{152}z_{153}$\\[5pt]
$Q_{132}=$ & $-100\,-20z_{12}\,-z_{12}^{2}\,-2z_{12}z_{13}\,-2z_{12}z_{14}\,-2z_{12}z_{15}$\\
 & $-2z_{12}z_{16}\,-2z_{12}z_{17}\,-2z_{12}z_{18}\,-2z_{12}z_{19}\,-2z_{12}z_{20}\,-20z_{13}$\\
 & $-z_{13}^{2}\,-2z_{13}z_{14}\,-2z_{13}z_{15}\,-2z_{13}z_{16}\,-2z_{13}z_{17}\,-2z_{13}z_{18}$\\
 & $-2z_{13}z_{19}\,-2z_{13}z_{20}\,-20z_{14}\,-z_{14}^{2}\,-2z_{14}z_{15}\,-2z_{14}z_{16}$\\
 & $-2z_{14}z_{17}\,-2z_{14}z_{18}\,-2z_{14}z_{19}\,-2z_{14}z_{20}\,-20z_{15}\,-z_{15}^{2}$\\
 & $-2z_{15}z_{16}\,-2z_{15}z_{17}\,-2z_{15}z_{18}\,-2z_{15}z_{19}\,-2z_{15}z_{20}\,-20z_{16}$\\
 & $-z_{16}^{2}\,-2z_{16}z_{17}\,-2z_{16}z_{18}\,-2z_{16}z_{19}\,-2z_{16}z_{20}\,-20z_{17}$\\
 & $-z_{17}^{2}\,-2z_{17}z_{18}\,-2z_{17}z_{19}\,-2z_{17}z_{20}\,-20z_{18}\,-z_{18}^{2}$\\
 & $-2z_{18}z_{19}\,-2z_{18}z_{20}\,-20z_{19}\,-z_{19}^{2}\,-2z_{19}z_{20}\,-20z_{20}$\\
 & $-z_{20}^{2}\,+z_{155}$\\[5pt]
$Q_{133}=$ & $-1\,+2z_{12}z_{123}\,+2z_{12}z_{124}\,+2z_{13}z_{123}\,+2z_{13}z_{124}\,+2z_{14}z_{123}$\\
 & $+2z_{14}z_{124}\,+2z_{15}z_{123}\,+2z_{15}z_{124}\,+2z_{16}z_{123}\,+2z_{16}z_{124}\,+2z_{17}z_{123}$\\
 & $+2z_{17}z_{124}\,+2z_{18}z_{123}\,+2z_{18}z_{124}\,+2z_{19}z_{123}\,+2z_{19}z_{124}\,+2z_{20}z_{123}$\\
 & $+2z_{20}z_{124}\,+20z_{123}\,+20z_{124}\,-z_{154}\,-z_{155}$\\[5pt]
$Q_{134}=$ & $z_{12}z_{150}\,+z_{13}z_{150}\,+z_{14}z_{150}\,+z_{15}z_{150}\,+z_{16}z_{150}\,+z_{17}z_{150}$\\
 & $+z_{18}z_{150}\,+z_{19}z_{150}\,+z_{20}z_{150}\,-z_{119}\,+z_{149}\,+10z_{150}$\\
 & $-z_{151}\,-z_{154}z_{156}$\\[5pt]
$Q_{135}=$ & $-10\,-6z_{4}\,-z_{4}^{2}\,+2z_{4}z_{123}\,+2z_{4}z_{124}\,-2z_{4}z_{173}$\\
 & $+6z_{123}\,+2z_{123}z_{173}\,+6z_{124}\,+2z_{124}z_{173}\,-z_{157}\,-6z_{173}$\\
 & $-z_{173}^{2}$\\[5pt]
$Q_{136}=$ & $-1\,+z_{4}z_{145}\,+z_{144}\,+3z_{145}\,+z_{145}z_{173}\,-z_{146}$\\
 & $-z_{157}z_{158}\,-z_{171}\,-z_{172}$\\[5pt]
$Q_{137}=$ & $-z_{123}z_{160}\,-z_{124}z_{160}\,+z_{161}$\\[5pt]
$Q_{138}=$ & $-1\,+z_{159}^{2}\,+z_{160}^{2}\,-z_{161}^{2}$\\[5pt]
$Q_{139}=$ & $-1\,+z_{4}z_{160}\,-z_{157}z_{162}\,+z_{159}\,+3z_{160}\,+z_{160}z_{173}$\\
 & $-z_{161}\,-z_{203}\,-z_{204}$\\[5pt]
$Q_{140}=$ & $-z_{123}z_{164}\,-z_{124}z_{164}\,+z_{165}$\\[5pt]
$Q_{141}=$ & $-1\,+z_{163}^{2}\,+z_{164}^{2}\,-z_{165}^{2}$\\[5pt]
$Q_{142}=$ & $-1\,+z_{4}z_{164}\,-z_{157}z_{166}\,+z_{163}\,+3z_{164}\,+z_{164}z_{173}$\\
 & $-z_{165}\,-z_{216}\,-z_{217}$\\[5pt]
$Q_{143}=$ & $-z_{123}z_{168}\,-z_{124}z_{168}\,+z_{169}$\\[5pt]
$Q_{144}=$ & $-1\,+z_{167}^{2}\,+z_{168}^{2}\,-z_{169}^{2}$\\[5pt]
$Q_{145}=$ & $-2\,-z_{4}\,+z_{4}z_{168}\,-z_{4}z_{171}\,-z_{4}z_{172}\,+z_{167}$\\
 & $+4z_{168}\,+z_{168}z_{173}\,-z_{169}\,-z_{170}\,-z_{170}z_{178}\,-z_{170}z_{179}$\\
 & $-3z_{171}\,-2z_{172}$\\[5pt]
$Q_{146}=$ & $-z_{123}z_{175}\,-z_{124}z_{175}\,+z_{176}$\\[5pt]
$Q_{147}=$ & $-1\,+z_{174}^{2}\,+z_{175}^{2}\,-z_{176}^{2}$\\[5pt]
$Q_{148}=$ & $-z_{21}\,-z_{21}z_{178}\,-z_{21}z_{179}\,+z_{180}$\\[5pt]
$Q_{149}=$ & $z_{4}z_{175}\,+z_{173}z_{175}\,+z_{174}\,+4z_{175}\,-z_{176}\,-z_{177}z_{180}$\\
 & $-z_{178}$\\[5pt]
$Q_{150}=$ & $-z_{123}z_{182}\,-z_{124}z_{182}\,+z_{183}$\\[5pt]
$Q_{151}=$ & $-1\,+z_{181}^{2}\,+z_{182}^{2}\,-z_{183}^{2}$\\[5pt]
$Q_{152}=$ & $z_{4}z_{182}\,+z_{173}z_{182}\,-z_{180}z_{184}\,+z_{181}\,+4z_{182}\,-z_{183}$\\
 & $-z_{185}$\\[5pt]
$Q_{153}=$ & $z_{178}\,+z_{179}\,-z_{185}\,-z_{186}$\\[5pt]
$Q_{154}=$ & $-z_{123}z_{188}\,-z_{124}z_{188}\,+z_{189}$\\[5pt]
$Q_{155}=$ & $-1\,+z_{187}^{2}\,+z_{188}^{2}\,-z_{189}^{2}$\\[5pt]
$Q_{156}=$ & $z_{4}z_{188}\,+z_{173}z_{188}\,-z_{180}z_{190}\,+z_{187}\,+4z_{188}\,-z_{189}$\\
 & $-z_{191}$\\[5pt]
$Q_{157}=$ & $z_{178}\,+z_{179}\,-z_{191}\,-z_{192}$\\[5pt]
$Q_{158}=$ & $-z_{123}z_{194}\,-z_{124}z_{194}\,+z_{195}$\\[5pt]
$Q_{159}=$ & $-1\,+z_{193}^{2}\,+z_{194}^{2}\,-z_{195}^{2}$\\[5pt]
$Q_{160}=$ & $z_{4}z_{194}\,+z_{173}z_{194}\,-z_{180}z_{196}\,+z_{193}\,+4z_{194}\,-z_{195}$\\
 & $-z_{197}$\\[5pt]
$Q_{161}=$ & $z_{178}\,+z_{179}\,-z_{197}\,-z_{198}$\\[5pt]
$Q_{162}=$ & $-z_{123}z_{200}\,-z_{124}z_{200}\,+z_{201}$\\[5pt]
$Q_{163}=$ & $-1\,+z_{199}^{2}\,+z_{200}^{2}\,-z_{201}^{2}$\\[5pt]
$Q_{164}=$ & $-z_{21}\,-z_{21}z_{203}\,-z_{21}z_{204}\,+z_{205}$\\[5pt]
$Q_{165}=$ & $z_{4}z_{200}\,+z_{173}z_{200}\,+z_{199}\,+4z_{200}\,-z_{201}\,-z_{202}z_{205}$\\
 & $-z_{203}$\\[5pt]
$Q_{166}=$ & $-z_{123}z_{207}\,-z_{124}z_{207}\,+z_{208}$\\[5pt]
$Q_{167}=$ & $-1\,+z_{206}^{2}\,+z_{207}^{2}\,-z_{208}^{2}$\\[5pt]
$Q_{168}=$ & $z_{4}z_{207}\,+z_{173}z_{207}\,-z_{205}z_{209}\,+z_{206}\,+4z_{207}\,-z_{208}$\\
 & $-z_{210}$\\[5pt]
$Q_{169}=$ & $z_{203}\,+z_{204}\,-z_{210}\,-z_{211}$\\[5pt]
$Q_{170}=$ & $-z_{123}z_{213}\,-z_{124}z_{213}\,+z_{214}$\\[5pt]
$Q_{171}=$ & $-1\,+z_{212}^{2}\,+z_{213}^{2}\,-z_{214}^{2}$\\[5pt]
$Q_{172}=$ & $-z_{21}\,-z_{21}z_{216}\,-z_{21}z_{217}\,+z_{218}$\\[5pt]
$Q_{173}=$ & $z_{4}z_{213}\,+z_{173}z_{213}\,+z_{212}\,+4z_{213}\,-z_{214}\,-z_{215}z_{218}$\\
 & $-z_{216}$\\[5pt]
$Q_{174}=$ & $-z_{123}z_{220}\,-z_{124}z_{220}\,+z_{221}$\\[5pt]
$Q_{175}=$ & $-1\,+z_{219}^{2}\,+z_{220}^{2}\,-z_{221}^{2}$\\[5pt]
$Q_{176}=$ & $z_{4}z_{220}\,+z_{173}z_{220}\,-z_{218}z_{222}\,+z_{219}\,+4z_{220}\,-z_{221}$\\
 & $-z_{223}$\\[5pt]
$Q_{177}=$ & $z_{216}\,+z_{217}\,-z_{223}\,-z_{224}$\\[5pt]
$Q_{178}=$ & $-z_{123}z_{226}\,-z_{124}z_{226}\,+z_{227}$\\[5pt]
$Q_{179}=$ & $-1\,+z_{225}^{2}\,+z_{226}^{2}\,-z_{227}^{2}$\\[5pt]
$Q_{180}=$ & $-z_{21}\,-z_{21}z_{171}\,-z_{21}z_{172}\,+z_{231}$\\[5pt]
$Q_{181}=$ & $z_{4}z_{226}\,+z_{173}z_{226}\,+z_{225}\,+4z_{226}\,-z_{227}\,-z_{228}z_{231}$\\
 & $-z_{229}$\\[5pt]
$Q_{182}=$ & $z_{171}\,+z_{172}\,-z_{229}\,-z_{230}$\\[5pt]
$Q_{183}=$ & $-z_{123}z_{233}\,-z_{124}z_{233}\,+z_{234}$\\[5pt]
$Q_{184}=$ & $-1\,+z_{232}^{2}\,+z_{233}^{2}\,-z_{234}^{2}$\\[5pt]
$Q_{185}=$ & $z_{4}z_{233}\,+z_{173}z_{233}\,-z_{180}z_{235}\,+z_{232}\,+4z_{233}\,-z_{234}$\\
 & $-z_{236}$\\[5pt]
$Q_{186}=$ & $z_{178}\,+z_{179}\,-z_{236}\,-z_{237}$\\[5pt]
$Q_{187}=$ & $1152\,+78z_{0}\,+3z_{0}^{2}\,+2z_{0}z_{1}\,+4z_{0}z_{2}\,+2z_{0}z_{3}$\\
 & $+2z_{0}z_{9}\,-2z_{0}z_{145}\,-2z_{0}z_{150}\,-2z_{0}z_{168}\,+32z_{1}\,+2z_{1}^{2}$\\
 & $+4z_{1}z_{2}\,+2z_{1}z_{9}\,-2z_{1}z_{150}\,-2z_{1}z_{164}\,+64z_{2}\,+5z_{2}^{2}$\\
 & $+6z_{2}z_{9}\,-4z_{2}z_{150}\,-2z_{2}z_{160}\,+24z_{3}\,+z_{3}^{2}\,-2z_{3}z_{168}$\\
 & $+34z_{9}\,+2z_{9}^{2}\,-2z_{9}z_{150}\,-2z_{9}z_{160}\,+52z_{12}\,+2z_{12}^{2}$\\
 & $+26z_{12}z_{13}\,+8z_{12}z_{23}\,+4z_{12}z_{26}\,+2z_{12}z_{30}\,+2z_{12}z_{32}\,+2z_{12}z_{34}$\\
 & $+2z_{12}z_{36}\,+4z_{12}z_{40}\,+8z_{12}z_{42}\,+4z_{12}z_{43}\,+2z_{12}z_{46}\,+4z_{12}z_{68}$\\
 & $-2z_{12}z_{175}\,-2z_{12}z_{233}\,+652z_{13}\,+170z_{13}^{2}\,+104z_{13}z_{23}\,+52z_{13}z_{26}$\\
 & $+26z_{13}z_{30}\,+26z_{13}z_{32}\,+26z_{13}z_{34}\,+26z_{13}z_{36}\,+52z_{13}z_{40}\,+104z_{13}z_{42}$\\
 & $+52z_{13}z_{43}\,+26z_{13}z_{46}\,+52z_{13}z_{68}\,-2z_{13}z_{182}\,-26z_{13}z_{233}\,+2z_{14}$\\
 & $+z_{14}^{2}\,-2z_{14}z_{188}\,+2z_{15}\,+z_{15}^{2}\,-2z_{15}z_{194}\,+2z_{16}$\\
 & $+z_{16}^{2}\,-2z_{16}z_{200}\,+2z_{17}\,+z_{17}^{2}\,-2z_{17}z_{207}\,+2z_{18}$\\
 & $+z_{18}^{2}\,-2z_{18}z_{213}\,+2z_{19}\,+z_{19}^{2}\,-2z_{19}z_{220}\,+2z_{20}$\\
 & $+z_{20}^{2}\,-2z_{20}z_{226}\,+200z_{23}\,+16z_{23}^{2}\,+16z_{23}z_{26}\,+8z_{23}z_{30}$\\
 & $+8z_{23}z_{32}\,+8z_{23}z_{34}\,+8z_{23}z_{36}\,+16z_{23}z_{40}\,+32z_{23}z_{42}\,+16z_{23}z_{43}$\\
 & $+8z_{23}z_{46}\,+16z_{23}z_{68}\,-8z_{23}z_{233}\,+100z_{26}\,+4z_{26}^{2}\,+4z_{26}z_{30}$\\
 & $+4z_{26}z_{32}\,+4z_{26}z_{34}\,+4z_{26}z_{36}\,+8z_{26}z_{40}\,+16z_{26}z_{42}\,+8z_{26}z_{43}$\\
 & $+4z_{26}z_{46}\,+8z_{26}z_{68}\,-4z_{26}z_{233}\,+50z_{30}\,+z_{30}^{2}\,+2z_{30}z_{32}$\\
 & $+2z_{30}z_{34}\,+2z_{30}z_{36}\,+4z_{30}z_{40}\,+8z_{30}z_{42}\,+4z_{30}z_{43}\,+2z_{30}z_{46}$\\
 & $+4z_{30}z_{68}\,-2z_{30}z_{233}\,+50z_{32}\,+z_{32}^{2}\,+2z_{32}z_{34}\,+2z_{32}z_{36}$\\
 & $+4z_{32}z_{40}\,+8z_{32}z_{42}\,+4z_{32}z_{43}\,+2z_{32}z_{46}\,+4z_{32}z_{68}\,-2z_{32}z_{233}$\\
 & $+50z_{34}\,+z_{34}^{2}\,+2z_{34}z_{36}\,+4z_{34}z_{40}\,+8z_{34}z_{42}\,+4z_{34}z_{43}$\\
 & $+2z_{34}z_{46}\,+4z_{34}z_{68}\,-2z_{34}z_{233}\,+50z_{36}\,+z_{36}^{2}\,+4z_{36}z_{40}$\\
 & $+8z_{36}z_{42}\,+4z_{36}z_{43}\,+2z_{36}z_{46}\,+4z_{36}z_{68}\,-2z_{36}z_{233}\,+100z_{40}$\\
 & $+4z_{40}^{2}\,+16z_{40}z_{42}\,+8z_{40}z_{43}\,+4z_{40}z_{46}\,+8z_{40}z_{68}\,-4z_{40}z_{233}$\\
 & $+200z_{42}\,+16z_{42}^{2}\,+16z_{42}z_{43}\,+8z_{42}z_{46}\,+16z_{42}z_{68}\,-8z_{42}z_{233}$\\
 & $+100z_{43}\,+4z_{43}^{2}\,+4z_{43}z_{46}\,+8z_{43}z_{68}\,-4z_{43}z_{233}\,+50z_{46}$\\
 & $+z_{46}^{2}\,+4z_{46}z_{68}\,-2z_{46}z_{233}\,+100z_{68}\,+4z_{68}^{2}\,-4z_{68}z_{233}$\\
 & $-z_{123}z_{239}\,-z_{124}z_{239}\,-24z_{145}\,+z_{145}^{2}\,-30z_{150}\,+z_{150}^{2}$\\
 & $-4z_{160}\,+z_{160}^{2}\,-2z_{164}\,+z_{164}^{2}\,-24z_{168}\,+z_{168}^{2}$\\
 & $-2z_{175}\,+z_{175}^{2}\,-2z_{182}\,+z_{182}^{2}\,-2z_{188}\,+z_{188}^{2}$\\
 & $-2z_{194}\,+z_{194}^{2}\,-2z_{200}\,+z_{200}^{2}\,-2z_{207}\,+z_{207}^{2}$\\
 & $-2z_{213}\,+z_{213}^{2}\,-2z_{220}\,+z_{220}^{2}\,-2z_{226}\,+z_{226}^{2}$\\
 & $-50z_{233}\,+z_{233}^{2}\,+z_{239}$\\[5pt]
$Q_{188}=$ & $-96\,-3z_{0}\,-2z_{1}\,-3z_{2}\,-2z_{3}\,-2z_{4}$\\
 & $-z_{5}\,-2z_{9}\,-3z_{12}\,-15z_{13}\,-2z_{14}\,-2z_{15}$\\
 & $-2z_{16}\,-2z_{17}\,-2z_{18}\,-2z_{19}\,-2z_{20}\,-z_{21}$\\
 & $-4z_{23}\,-2z_{26}\,-z_{30}\,-z_{32}\,-z_{34}\,-z_{36}$\\
 & $-2z_{40}\,-4z_{42}\,-2z_{43}\,-z_{46}\,-2z_{68}\,-z_{119}$\\
 & $-z_{171}\,-z_{172}\,-z_{173}\,-z_{203}\,-z_{204}\,+z_{240}$\\[5pt]
\bottomrule\end{longtable}
\endgroup


\section{Coordinate meanings}

These names identify compiler coordinates; they impose no extra conditions
on the variables.  The preceding residual list is the complete arithmetic
definition.  Names of eliminated variables occur only in the supplied
restoration certificates.
\begingroup\footnotesize
\begin{longtable}{@{}r p{0.40\linewidth}r p{0.40\linewidth}@{}}
\toprule Index & Name & Index & Name\\\midrule\endhead
0 & \pathfile{t0} & 121 & \pathfile{mul.426}\\
1 & \pathfile{x} & 122 & \pathfile{mul.427}\\
2 & \pathfile{fact0.output} & 123 & \pathfile{shared.lower_A}\\
3 & \pathfile{mul.1} & 124 & \pathfile{shared.A_slack}\\
4 & \pathfile{reciprocal.Bminus2} & 125 & \pathfile{shared.Zminus1}\\
5 & \pathfile{pow0.output} & 126 & \pathfile{mul.431}\\
6 & \pathfile{fact0.remainder} & 127 & \pathfile{mul.432}\\
7 & \pathfile{fact0.slack} & 128 & \pathfile{modular.root_quotient}\\
8 & \pathfile{endpoint.u0} & 129 & \pathfile{index_cap.X}\\
9 & \pathfile{endpoint.r0} & 130 & \pathfile{index_cap.U}\\
10 & \pathfile{endpoint.ut} & 131 & \pathfile{index_cap.V}\\
11 & \pathfile{mul.12} & 132 & \pathfile{index_cap.S}\\
12 & \pathfile{crt.encoded.a} & 133 & \pathfile{index_cap.R}\\
13 & \pathfile{crt.encoded.b} & 134 & \pathfile{index_cap.q4minus1}\\
14 & \pathfile{crt.encoded.c} & 135 & \pathfile{mul.441}\\
15 & \pathfile{crt.encoded.d} & 136 & \pathfile{mul.443}\\
16 & \pathfile{crt.encoded.S} & 137 & \pathfile{mul.444}\\
17 & \pathfile{crt.encoded.Sp} & 138 & \pathfile{mul.445}\\
18 & \pathfile{crt.encoded.u} & 139 & \pathfile{index_cap.Vquotient_minus1}\\
19 & \pathfile{crt.encoded.vq} & 140 & \pathfile{mul.446}\\
20 & \pathfile{crt.encoded_index_minus_one} & 141 & \pathfile{index_cap.Bquotient}\\
21 & \pathfile{pow1.output} & 142 & \pathfile{index_cap.Squotient}\\
22 & \pathfile{crt.index_quotient} & 143 & \pathfile{index_cap.Rquotient}\\
23 & \pathfile{mul.41} & 144 & \pathfile{index0.X}\\
24 & \pathfile{mul.48} & 145 & \pathfile{index0.Y}\\
25 & \pathfile{mul.53} & 146 & \pathfile{mul.455}\\
26 & \pathfile{mul.60} & 147 & \pathfile{base0.T}\\
27 & \pathfile{mul.65} & 148 & \pathfile{pow0.quotient}\\
28 & \pathfile{mul.79} & 149 & \pathfile{index1.X}\\
29 & \pathfile{add.89} & 150 & \pathfile{index1.Y}\\
30 & \pathfile{add.94} & 151 & \pathfile{mul.465}\\
31 & \pathfile{add.106} & 152 & \pathfile{base1.T}\\
32 & \pathfile{add.111} & 153 & \pathfile{pow1.quotient}\\
33 & \pathfile{add.126} & 154 & \pathfile{base2.T}\\
34 & \pathfile{add.131} & 155 & \pathfile{mul.476}\\
35 & \pathfile{add.140} & 156 & \pathfile{pow2.quotient}\\
36 & \pathfile{add.145} & 157 & \pathfile{base3.T}\\
37 & \pathfile{mul.148} & 158 & \pathfile{pow3.quotient}\\
38 & \pathfile{mul.153} & 159 & \pathfile{index2.X}\\
39 & \pathfile{mul.160} & 160 & \pathfile{index2.Y}\\
40 & \pathfile{mul.165} & 161 & \pathfile{mul.493}\\
41 & \pathfile{add.176} & 162 & \pathfile{pow4.quotient}\\
42 & \pathfile{add.181} & 163 & \pathfile{index3.X}\\
43 & \pathfile{mul.184} & 164 & \pathfile{index3.Y}\\
44 & \pathfile{mul.190} & 165 & \pathfile{mul.499}\\
45 & \pathfile{add.204} & 166 & \pathfile{pow5.quotient}\\
46 & \pathfile{add.209} & 167 & \pathfile{index4.X}\\
47 & \pathfile{add.210} & 168 & \pathfile{index4.Y}\\
48 & \pathfile{mul.219} & 169 & \pathfile{mul.505}\\
49 & \pathfile{mul.224} & 170 & \pathfile{choose0.modular_quotient}\\
50 & \pathfile{mul.230} & 171 & \pathfile{choose0.modular_remainder}\\
51 & \pathfile{mul.237} & 172 & \pathfile{choose0.modular_rslack}\\
52 & \pathfile{mul.242} & 173 & \pathfile{choose0.modular_oslack}\\
53 & \pathfile{mul.248} & 174 & \pathfile{index5.X}\\
54 & \pathfile{add.253} & 175 & \pathfile{index5.Y}\\
55 & \pathfile{add.258} & 176 & \pathfile{mul.517}\\
56 & \pathfile{mul.261} & 177 & \pathfile{crt.range.a.modular_quotient}\\
57 & \pathfile{mul.268} & 178 & \pathfile{crt.range.a.modular_remainder}\\
58 & \pathfile{mul.273} & 179 & \pathfile{crt.range.a.modular_rslack}\\
59 & \pathfile{mul.279} & 180 & \pathfile{mul.518}\\
60 & \pathfile{mul.286} & 181 & \pathfile{index6.X}\\
61 & \pathfile{mul.291} & 182 & \pathfile{index6.Y}\\
62 & \pathfile{mul.304} & 183 & \pathfile{mul.526}\\
63 & \pathfile{add.307} & 184 & \pathfile{crt.range.b.modular_quotient}\\
64 & \pathfile{mul.309} & 185 & \pathfile{crt.range.b.modular_remainder}\\
65 & \pathfile{mul.322} & 186 & \pathfile{crt.range.b.modular_rslack}\\
66 & \pathfile{add.332} & 187 & \pathfile{index7.X}\\
67 & \pathfile{add.337} & 188 & \pathfile{index7.Y}\\
68 & \pathfile{mul.340} & 189 & \pathfile{mul.534}\\
69 & \pathfile{mul.343} & 190 & \pathfile{crt.range.c.modular_quotient}\\
70 & \pathfile{sqd.0} & 191 & \pathfile{crt.range.c.modular_remainder}\\
71 & \pathfile{mul.344} & 192 & \pathfile{crt.range.c.modular_rslack}\\
72 & \pathfile{sqd.1} & 193 & \pathfile{index8.X}\\
73 & \pathfile{mul.346} & 194 & \pathfile{index8.Y}\\
74 & \pathfile{sqd.2} & 195 & \pathfile{mul.542}\\
75 & \pathfile{mul.348} & 196 & \pathfile{crt.range.d.modular_quotient}\\
76 & \pathfile{sqd.3} & 197 & \pathfile{crt.range.d.modular_remainder}\\
77 & \pathfile{mul.350} & 198 & \pathfile{crt.range.d.modular_rslack}\\
78 & \pathfile{sqd.4} & 199 & \pathfile{index9.X}\\
79 & \pathfile{mul.352} & 200 & \pathfile{index9.Y}\\
80 & \pathfile{sqd.5} & 201 & \pathfile{mul.550}\\
81 & \pathfile{mul.354} & 202 & \pathfile{crt.range.S.modular_quotient}\\
82 & \pathfile{sqd.6} & 203 & \pathfile{crt.range.S.modular_remainder}\\
83 & \pathfile{mul.356} & 204 & \pathfile{crt.range.S.modular_rslack}\\
84 & \pathfile{sqd.7} & 205 & \pathfile{mul.551}\\
85 & \pathfile{mul.358} & 206 & \pathfile{index10.X}\\
86 & \pathfile{sqd.8} & 207 & \pathfile{index10.Y}\\
87 & \pathfile{sqd.9} & 208 & \pathfile{mul.559}\\
88 & \pathfile{sqd.10} & 209 & \pathfile{crt.range.Sp.modular_quotient}\\
89 & \pathfile{mul.360} & 210 & \pathfile{crt.range.Sp.modular_remainder}\\
90 & \pathfile{mul.362} & 211 & \pathfile{crt.range.Sp.modular_rslack}\\
91 & \pathfile{sqd.11} & 212 & \pathfile{index11.X}\\
92 & \pathfile{mul.363} & 213 & \pathfile{index11.Y}\\
93 & \pathfile{sqd.12} & 214 & \pathfile{mul.567}\\
94 & \pathfile{mul.366} & 215 & \pathfile{crt.range.u.modular_quotient}\\
95 & \pathfile{mul.368} & 216 & \pathfile{crt.range.u.modular_remainder}\\
96 & \pathfile{mul.369} & 217 & \pathfile{crt.range.u.modular_rslack}\\
97 & \pathfile{mul.371} & 218 & \pathfile{mul.568}\\
98 & \pathfile{mul.373} & 219 & \pathfile{index12.X}\\
99 & \pathfile{mul.374} & 220 & \pathfile{index12.Y}\\
100 & \pathfile{sqd.13} & 221 & \pathfile{mul.576}\\
101 & \pathfile{mul.375} & 222 & \pathfile{crt.range.vq.modular_quotient}\\
102 & \pathfile{mul.377} & 223 & \pathfile{crt.range.vq.modular_remainder}\\
103 & \pathfile{mul.378} & 224 & \pathfile{crt.range.vq.modular_rslack}\\
104 & \pathfile{sqd.14} & 225 & \pathfile{index13.X}\\
105 & \pathfile{mul.379} & 226 & \pathfile{index13.Y}\\
106 & \pathfile{mul.380} & 227 & \pathfile{mul.584}\\
107 & \pathfile{mul.381} & 228 & \pathfile{crt.range.index.modular_quotient}\\
108 & \pathfile{mul.382} & 229 & \pathfile{crt.range.index.modular_remainder}\\
109 & \pathfile{mul.383} & 230 & \pathfile{crt.range.index.modular_rslack}\\
110 & \pathfile{mul.384} & 231 & \pathfile{mul.585}\\
111 & \pathfile{mul.385} & 232 & \pathfile{index14.X}\\
112 & \pathfile{mul.386} & 233 & \pathfile{index14.Y}\\
113 & \pathfile{mul.387} & 234 & \pathfile{mul.593}\\
114 & \pathfile{mul.388} & 235 & \pathfile{crt.range.derived_label_sum.modular_quotient}\\
115 & \pathfile{mul.391} & 236 & \pathfile{crt.range.derived_label_sum.modular_remainder}\\
116 & \pathfile{mul.393} & 237 & \pathfile{crt.range.derived_label_sum.modular_rslack}\\
117 & \pathfile{mul.396} & 238 & \pathfile{index_cap.aggregate_slack}\\
118 & \pathfile{crt.nonnegative_quotient} & 239 & \pathfile{index_cap.error_square_quotient}\\
119 & \pathfile{pow2.output} & 240 & \pathfile{add.639}\\
120 & \pathfile{mul.425} &  & \\
\bottomrule\end{longtable}
\endgroup


\section{Reproduction and dependency map}

The top-level \pathfile{README.md} gives the selected files and exact
replay command.  The direct subdirectory contains the frozen compiler
inputs, generator, whole expanded quartics, and elimination certificates.
The positive subdirectory contains its fixed coefficient payloads,
source equations, literal schedules, ordinary proof, and independent
arithmetic checks.  Retained upstream licenses and notices apply to
their corresponding files.

The release verifier checks the file manifest, reproduces the selected
mathematical artifacts, checks the complete principal residual appendix,
and writes its receipt to a separate location.  Source and artifact
hashes pin the comparison and prevent an accidental mixture of constants
or compiler versions.  Hashes authenticate finite bytes; they do not
replace the unbounded mathematical arguments.

\begingroup\small
\begin{thebibliography}{9}
\setlength{\itemsep}{2pt}\setlength{\parskip}{0pt}
\bibitem{repo} ProveIt Contributors, \emph{ProveIt}, snapshot
\pathfile{f4d8db99e060339f3306a04d507c35bc1d0d68c5}, 24 October 2026.
\url{https://github.com/VladimirReshetnikov/ProveIt}.
\bibitem{binomial} ProveIt incoming research packet,
\emph{Shared Binomial Certificates and Short Carriers for Con(ZFC)},
24 October 2026.  Archive
\pathfile{ConZFC_Shared_Binomial_Certificates_2026-10-24.zip},
\pathfile{docs/incoming} at the pinned snapshot.  Its direct ordinary
proof and finite source are retained in this packet.
\bibitem{cap261} ProveIt incoming research packet,
\emph{One Pell Bound and One Error Congruence}, 24 October 2026.
Archive \pathfile{ConZFC_Shared_Pell_Bound_261_2026-10-24.zip},
same pinned directory.  The one-anchor theorem is retained as
\pathfile{CAP_THEOREM.md}.
\bibitem{gapfusion} ProveIt Contributors,
\emph{Coprime Fusion with the Auxiliary-Index Gap Majorant},
\pathfile{Computability/DiophantineConsistency/candidates/ten-positive-gap-fusion/PROOF.md},
same pinned snapshot.
\bibitem{unified} ProveIt incoming research packet,
\emph{Round 15 Unified Binomial and Even Pell Index Equations},
24 October 2026. Archive
\pathfile{ConZFC_Unified_Binomial_Even_Indices_Round15_2026-10-24.zip},
same pinned directory. The raw ceiling bound is retained in
\pathfile{direct/source/UNIFIED_PRODUCT_PROOF.md}.
\bibitem{bayerdavid} J. Bayer and M. David,
\emph{A Formal Proof of Complexity Bounds on Diophantine Equations},
ITP 2025, LIPIcs 352, Article 3, 2025.
\url{https://doi.org/10.4230/LIPIcs.ITP.2025.3}.
\bibitem{bayeretal} J. Bayer, M. David, M. Ha{\ss}ler, Y. Matiyasevich,
and D. Schleicher, \emph{Diophantine Equations over $\Z$: Universal Bounds
and Parallel Formalization}, arXiv:2506.20909v3, 22 October 2026.
\url{https://arxiv.org/abs/2506.20909}.
\end{thebibliography}
\endgroup
\end{document}
